\chapter{Introduction to von Neumann algebras}

Chapter~1 ended with a proposal. A subsystem should be defined not by splitting the Hilbert space, but by
the collection of operators that an observer confined to the subsystem can use. This chapter makes that
proposal precise. By the end of it you will know what kind of mathematical object $\M$ must be, how such
objects are sorted into types, and how a Hilbert space can be \emph{built} from an algebra and a state
instead of being assumed at the start. Most definitions below come with a worked example that uses actual
numbers, often small matrices you can multiply out by hand. The definitions are abstract on first reading,
and watching them act on something small is the quickest way to make them concrete.

\section{Systems with an infinite amount of entanglement (recap)}

Chapter~1 gave three examples: the chain of $N$ Bell pairs, lattice gauge theory, and a quantum field theory
cut in half. They motivate everything in this chapter. The lesson drawn from them was this: a subsystem
should be defined by \emph{which operators are available to act with}, not by how the Hilbert space happens
to split. This chapter turns that idea into a mathematical definition.

\section{Von Neumann algebras as subsystems}

\subsection*{Setting the stage: what a Hilbert space and an operator are, restated carefully}

The new definition is a statement \emph{about} Hilbert spaces and operators. So we first state exactly what
those are. Vague pictures of ``vectors'' and ``operators'' are not precise enough for what follows.

A Hilbert space $\HH$ is a vector space over the complex numbers with two extra properties. First, it has an
inner product $\braket{\cdot|\cdot}$. This is a rule that assigns a complex number $\braket{\xi|\eta}$ to each
pair of vectors. It is linear in the second argument and conjugate-linear in the first, and it satisfies
$\braket{\xi|\xi}\ge0$, with equality only for the zero vector. Second, it is complete. A \emph{Cauchy
sequence} is a sequence of vectors whose terms get arbitrarily close to each other, so that it ``ought to''
converge. Completeness says that every Cauchy sequence really does converge to some vector in $\HH$. For a
spin-$\tfrac12$ particle, $\HH=\mathbb C^2$. Its vectors are pairs of complex numbers $(c_1,c_2)$, usually
written $c_1\ket0+c_2\ket1$. The inner product is the dot product with complex conjugation,
$\braket{\xi|\eta}=\xi_1^*\eta_1+\xi_2^*\eta_2$. A von Neumann algebra needs nothing more exotic than this.
The infinite-dimensional case (a quantum field, say) uses exactly the same definitions, with infinitely many
components instead of two.

An operator $A$ on $\HH$ is a rule that turns vectors into vectors, $A:\HH\to\HH$, and is linear:
$A(c_1\ket{\xi_1}+c_2\ket{\xi_2})=c_1 A\ket{\xi_1}+c_2A\ket{\xi_2}$. On $\mathbb C^2$, every linear operator is
a $2\times2$ matrix, and $A\ket\xi$ is ordinary matrix-vector multiplication. The \textbf{Hermitian
conjugate} (or adjoint) $A^\dagger$ of $A$ is the unique operator with
$\braket{\xi|A\eta}=\braket{A^\dagger\xi|\eta}$ for all $\xi,\eta$. For a finite matrix, this is the familiar
rule ``transpose and complex-conjugate every entry.'' $A$ is \textbf{self-adjoint} (or Hermitian) if
$A=A^\dagger$. Self-adjoint operators are the ones that represent physical observables, because their
eigenvalues are real numbers, the kind of number a measurement can return. $B(\HH)$ denotes the set of
\emph{all} bounded operators on $\HH$. Boundedness was defined in Chapter~1. Informally, a bounded operator
does not stretch the length of any vector by more than some fixed factor.

\subsection*{The new definition of subsystem}

Here is the definition, with every piece spelled out. Suppose you have access to only some subset of
operators, $\M\subset B(\HH)$. This is not everything. It is only what your apparatus, in your location, with
your capabilities, can measure or apply. Let $\ket\Psi\in\HH$ be a state of the full system. Using only $\M$,
you can do two things:
\begin{itemize}
\item compute expectation values $\braket{\Psi|A|\Psi}$ for $A\in\M$, which predict the average reading of a
measurement of $A$ over many repetitions;
\item act on the state, producing $B\ket\Psi$ for $B\in\M$, which changes the system by an operation you have
access to.
\end{itemize}
Whatever $\M$ is, it should be closed under three operations, because an observer with access to $\M$ can
already do each of them. If you can measure or apply $A$ and $B$, you can measure or apply $A+B$ (do both and
add the results) and $AB$ (do one after the other). You can also multiply by complex numbers. And if $A$ is
available, so is $A^\dagger$. For a self-adjoint $A$ this is automatic, since $A^\dagger=A$. In general it
says that your toolkit is closed under taking adjoints. A set closed under sums, products, and adjoints is
called a \textbf{$*$-subalgebra} of $B(\HH)$. The $*$ refers to the adjoint, which mathematics texts often
write as $A^*$ instead of $A^\dagger$. In these notes every $*$-subalgebra is also taken to contain the identity
operator $\id$, which is the operation ``do nothing.''

This alone is not quite enough. We also want $\M$ to be \emph{complete}: if some operator can be approximated
arbitrarily well by elements of $\M$, it should count as an element of $\M$ too. Otherwise the set of things
you have access to would have artificial gaps. Here the story becomes subtle. There are two different,
inequivalent ways to make ``approximated arbitrarily well'' precise, and the choice between them matters a
great deal. The next section explains both.

With this notion of subsystem in hand, the \textbf{complement} of $\M$ is defined to be its commutant,
\[
\M' = \{B\in B(\HH) : BA=AB\ \ \forall A\in\M\} .
\]
One misreading of this definition is common enough to address right away. $\M'$ is \emph{not} the set of
operators that commute \emph{with one another}. It is the set of operators that \emph{each} commute with
\emph{every} element of $\M$. Two operators inside $\M'$ can fail to commute with each other. In fact, when
$\M'$ is a rich enough algebra (the usual case), it contains non-commuting pairs, just as $\M$ does. What
$\M'$ guarantees is only this: acting with anything in $\M'$ produces no interference with any measurement
made using $\M$. That is the operational meaning of ``complement of a subsystem'' here. It is not ``the other
half of a tensor factorization you could point to.'' It is ``everything that, as a matter of algebra, cannot
disturb what $\M$ measures.'' A fuller treatment, with the functional-analysis machinery spelled out, is given
in W.~Thirring's \emph{Quantum Mathematical Physics}. Nothing beyond what is given here is needed for the rest
of these notes.

\begin{physconnect}[: Spacelike Commutativity]
You already know one instance of this idea from relativistic field theory. It is the cleanest bridge from
the abstract definition to something you can compute. For a free scalar field $\phi$, the commutator of the
field at two spacetime points is a fixed function times the identity operator, so it does not depend on the
state:
\[
[\phi(x),\phi(y)] = i\Delta(x-y) .
\]
Here $\Delta$, the Pauli--Jordan function, is given by an integral over the on-shell momenta of the field. A
standard textbook fact is that $\Delta(x-y)=0$ whenever $x-y$ is spacelike. The reason uses two properties of
$\Delta$. First, $\Delta$ is Lorentz invariant: $\Delta(\Lambda z)=\Delta(z)$ for every proper orthochronous
Lorentz transformation $\Lambda$. Second, $\Delta$ is odd, $\Delta(-z)=-\Delta(z)$, because the commutator is
antisymmetric, $[\phi(x),\phi(y)]=-[\phi(y),\phi(x)]$. Now let $z$ be a spacelike vector in $3+1$ dimensions.
Then there is a proper orthochronous Lorentz transformation with $\Lambda z=-z$. (Boost to a frame in which
$z$ has no time component, rotate by $180^\circ$ about an axis perpendicular to $z$, and boost back.) Hence
$\Delta(z)=\Delta(-z)=-\Delta(z)$, so $\Delta(z)=0$. For a timelike $z$ no such $\Lambda$ exists, because these
transformations never exchange past and future. This is why microcausality, the statement that field
operators at spacelike separation commute, comes out of the free-field construction automatically instead
of being imposed as an extra axiom.

Now translate this into the language of this chapter. Let $\M(O)$ be the algebra generated by field
operators smeared over a spacetime region $O$. Let $O'$ be its causal complement, the set of points
spacelike separated from all of $O$. The vanishing of $\Delta$ at spacelike separation says exactly that
$\M(O')\subseteq\M(O)'$. This is the \emph{locality} axiom for the local algebras of a relativistic field
theory, stated in general in Chapter~4. So in a relativistic field theory, every operator localized
spacelike to $O$ lies in the commutant $\M(O)'$.

The abstract definition of $\M'$ adds one thing to this familiar fact. The commutant is defined by a purely
algebraic condition, ``commutes with everything in $\M$,'' which never mentions spacetime. Nothing in the
definition forces every such operator to be localized in $O'$. In general, $\M(O)'$ can be strictly larger
than $\M(O')$. When the two are equal, the theory is said to satisfy \emph{Haag duality} for the region $O$
(Chapter~4). Haag duality holds in many standard cases but not in all, and whether it holds depends on the
theory, the region, and the representation. So ``commutes with $\M$'' is a more general notion than ``is
spacelike separated from $O$.'' Chapter~8 uses exactly this extra room: there, bulk causal structure is read
off from boundary commutants that are larger than boundary causality alone would require.
\end{physconnect}

\section{Basic properties: \texorpdfstring{$C^*$}{C*} vs.\ von Neumann algebras, and the double commutant theorem}

\subsection*{Two notions of ``the limit of a sequence of operators''}

Take a sequence of operators $A_1,A_2,A_3,\dots$ in $\M$, and ask whether it converges to some operator $A$.
For ordinary numbers, ``converges'' has one obvious meaning. For operators there are several natural
meanings, and they are genuinely different.

\textbf{Norm convergence.} Every bounded operator has a norm, $\|A\|$. It is the smallest number such that
$\|A\ket\psi\|\le\|A\|\,\|\ket\psi\|$ for every vector $\ket\psi$: the ``maximum stretch factor'' of
Chapter~1. The norm satisfies the identity
\begin{equation}
\|A^\dagger A\|=\|A\|^2 .
\label{eq:Cstar-identity}
\end{equation}
It is easy to check on a simple case. Take $A=\begin{psmallmatrix}0&1\\0&0\end{psmallmatrix}$ acting on
$\mathbb C^2$. Then $A^\dagger A = \begin{psmallmatrix}0&0\\0&1\end{psmallmatrix}$, whose largest eigenvalue
is $1$, so $\|A^\dagger A\|=1$. The norm $\|A\|$ is the largest singular value of $A$, which is also $1$. So
$\|A^\dagger A\|=1=1^2=\|A\|^2$, as the identity requires. A sequence $A_n$ \textbf{converges in norm} to $A$
if $\|A_n-A\|\to0$. This says that the largest possible discrepancy between $A_n\ket\psi$ and $A\ket\psi$,
taken over \emph{all} unit vectors $\ket\psi$ at once, shrinks to zero. It is a strong, uniform statement.

\textbf{Strong convergence.} $A_n$ \textbf{converges strongly} to $A$ if $\|(A_n-A)\ket\psi\|\to0$ for each
fixed vector $\ket\psi$. Each vector is tested separately, so the rate of convergence may differ from one
vector to another.

\textbf{Weak convergence.} $A_n$ \textbf{converges weakly} to $A$ if $\braket{\xi|A_n|\eta}\to
\braket{\xi|A|\eta}$ for every fixed pair of vectors $\ket\xi,\ket\eta$. This asks only that individual
matrix elements converge, one pair of vectors at a time. It is the weakest of the three demands.

Each notion implies the next one. If $\|A_n-A\|\to0$, then $\|(A_n-A)\ket\psi\|\le\|A_n-A\|\,\|\ket\psi\|\to0$,
so norm convergence implies strong convergence. And $|\braket{\xi|(A_n-A)|\eta}|\le\|\ket\xi\|\,
\|(A_n-A)\ket\eta\|$ by the Cauchy--Schwarz inequality, so strong convergence implies weak convergence. In
finite dimensions the three notions coincide completely. This is why the distinction never comes up in an
ordinary quantum mechanics course. In infinite dimensions they differ, and the gap between them is where
infinite-dimensional phenomena (quantum field theory, the $N\to\infty$ limit) live. The next example shows
the difference.

\begin{quote}
\textit{Worked example: weak vs.\ strong vs.\ norm convergence in $\ell^2(\mathbb N)$.} Let
$\HH = \ell^2(\mathbb N)$ be the infinite-dimensional Hilbert space of square-summable sequences, with
standard orthonormal basis $\{\ket 1, \ket 2, \ket 3, \dots\}$. Consider the sequence of rank-one projections
\[
P_n \equiv \ket n\bra n .
\]
We test the convergence of $P_n$ as $n\to\infty$ in each of the three senses.
\begin{enumerate}
\item \textbf{Weak convergence: $P_n \to 0$ weakly.} Take any two vectors
$\ket\xi = \sum_{k=1}^\infty c_k\ket k$ and $\ket\eta = \sum_{k=1}^\infty d_k\ket k$, with
$\sum |c_k|^2 < \infty$ and $\sum |d_k|^2 < \infty$. The matrix element is
\[
\braket{\xi|P_n|\eta} = \braket{\xi|n}\braket{n|\eta} = c_n^* d_n .
\]
The Cauchy--Schwarz inequality for series gives
$\sum_{n=1}^\infty |c_n d_n| \le \sqrt{\sum |c_n|^2}\sqrt{\sum |d_n|^2} < \infty$. The terms of a convergent
series tend to zero, so $c_n^* d_n\to0$. Hence $\braket{\xi|P_n|\eta} \to 0$ for \emph{every} pair of vectors,
and $P_n$ converges weakly to the zero operator.

\item \textbf{Strong convergence: $P_n \to 0$ strongly.} For a fixed vector $\ket\psi = \sum c_k\ket k$ we
compute
\[
\|P_n\ket\psi\|^2 = \|\ket n\braket{n|\psi}\|^2 = |c_n|^2 \to 0 \quad \text{as } n\to\infty .
\]
The limit is zero because $|c_n|^2$ are the terms of the convergent series $\sum_k|c_k|^2$. Hence $P_n$
converges strongly to $0$.

\item \textbf{Norm convergence: $P_n$ does not converge in norm.} The operator norm takes the largest stretch
over all unit vectors at once. The largest stretch is reached at $\ket\psi=\ket n$:
\[
\|P_n - 0\| = \sup_{\|\psi\|=1} \|P_n\ket\psi\| = \|P_n\ket n\| = \|\ket n\| = 1 .
\]
So $\|P_n - 0\| = 1$ for every $n$, and this never shrinks. Hence $P_n$ does \emph{not} converge to $0$ in
norm. It does not converge in norm to anything else either. For $n\ne m$ the operator $P_n-P_m$ has
eigenvalues $+1$ and $-1$, so $\|P_n-P_m\|=1$, and the sequence is not even a Cauchy sequence in norm.
\end{enumerate}
This is the standard example to keep in mind. Every matrix element settles down to zero, but the operator as
a whole never becomes small in norm.
\end{quote}

A $*$-subalgebra $\M$ that is complete under norm convergence is called a \textbf{$C^*$-algebra}. Here
``complete'' means that every norm-convergent sequence in $\M$ has its limit back in $\M$. A $*$-subalgebra
that is complete under weak convergence is called a \textbf{von Neumann algebra}. (Strictly speaking, the
weak closure is defined with nets rather than sequences. Sequences give the right intuition, and nothing in
these notes depends on the difference.) Every von Neumann algebra is automatically a $C^*$-algebra. To see
this, suppose $A_n\in\M$ and $A_n\to A$ in norm. Then $A_n\to A$ weakly as well, so if $\M$ is weakly closed,
$A\in\M$. Weak closure is the stronger requirement, because it asks $\M$ to contain the limits of more
sequences. So von Neumann algebras are the more complete, more restrictive notion.

Which notion is the right one for a physical subsystem? Weak convergence, for a simple reason. What a
laboratory measures is a matrix element $\braket{\xi|A|\eta}$: an expectation value, or more generally a
transition amplitude. Weak convergence says exactly that each of these measurable numbers settles down. So
von Neumann algebras, and not the less complete $C^*$-algebras, capture ``the set of operators whose physical
predictions are under control.'' That is why they are the main tool in the rest of these notes.

There is one more difference between the two notions, and it matters for everything that follows. A
$C^*$-algebra can be defined abstractly, with no Hilbert space anywhere. It is a vector space with a product,
an adjoint, and a norm that satisfies the identity~\eqref{eq:Cstar-identity}, and it is complete in that
norm. You could specify one by its multiplication table and its norm, much as you specify a finite group by
its multiplication table, without saying what vectors it acts on. The definition of a von Neumann algebra
given above cannot be stated this way. Its completeness condition, weak convergence, is phrased in terms of
vectors $\ket\xi,\ket\eta$ in a specific $\HH$. (Mathematicians do have an abstract characterization of von
Neumann algebras, due to Sakai, but it is less direct and is not needed here.) This gap is exactly what the
GNS construction, worked through later in this chapter, is designed to close. Starting from nothing but an
abstract $C^*$-algebra and a state on it, GNS builds the missing Hilbert space. Only then can you ask whether
the algebra, represented on that Hilbert space, is also weakly closed.

\subsection*{The double commutant theorem}

Checking weak closure directly means verifying that \emph{every} weakly convergent sequence in $\M$ has its
limit in $\M$. That sounds like infinitely many checks. Von Neumann's double commutant theorem replaces this
topological chore with an algebraic check. For a $*$-subalgebra $\M\subset B(\HH)$ that contains the
identity,
\[
\M = \M'' \iff \M \text{ is weakly closed (a von Neumann algebra).}
\]
Here $\M''\equiv(\M')'$ is the commutant of the commutant. Read the logic slowly. You take the complement of
your subsystem, $\M'$: everything that cannot disturb $\M$. Then you take the complement of \emph{that},
$\M''$: everything that cannot disturb anything in $\M'$. The theorem says that doing this twice brings you
back exactly to where you started, provided $\M$ was already a von Neumann algebra. If $\M$ is only a
$*$-subalgebra (closed under sums, products, and adjoints, but not weakly closed), then $\M''$ is in general
strictly larger than $\M$. It is the \emph{smallest} von Neumann algebra containing $\M$: $\M$ together with
every operator that can be approximated weakly by elements of $\M$. This is not an obvious fact. It is a real
theorem, and its proof uses the geometry of Hilbert space. Its value is that a purely algebraic operation
(take the commutant twice) carries out a topological completion (weak closure) for you, with no limits in
the definition at all.

\begin{quote}
\textit{Worked example, done completely, in four dimensions.} Let $\HH=\mathbb C^2\otimes\mathbb C^2$ — two
qubits, call them $R$ and $L$, so $\HH$ is 4-dimensional. Let $\M$ be every operator of the form $A\otimes
\id_L$ for $A$ an arbitrary $2\times2$ matrix (every operation you could perform using only apparatus touching
qubit $R$, doing nothing to qubit $L$). As a vector space, $\M$ is spanned by four matrices: the identity
$\id_2$ and the three Pauli matrices $\sigma_x,\sigma_y,\sigma_z$, each tensored with $\id_L$:
\[
\id_2=\begin{pmatrix}1&0\\0&1\end{pmatrix},\ \
\sigma_x=\begin{pmatrix}0&1\\1&0\end{pmatrix},\ \
\sigma_y=\begin{pmatrix}0&-i\\i&0\end{pmatrix},\ \
\sigma_z=\begin{pmatrix}1&0\\0&-1\end{pmatrix}.
\]
Every $2\times2$ matrix is some combination $c_0\id_2+c_1\sigma_x+c_2\sigma_y+c_3\sigma_z$ (this is just
saying the four Pauli matrices, including the identity, form a basis for all $2\times2$ matrices — a standard
fact, and easy to check directly: writing out a general $2\times2$ matrix $\begin{psmallmatrix}a&b\\c&d
\end{psmallmatrix}$ and matching coefficients gives $c_0=\tfrac{a+d}2$, $c_3=\tfrac{a-d}2$,
$c_1=\tfrac{b+c}2$, $c_2=\tfrac{b-c}{2i}$, which always has a solution).

To find $\M'$, look for every $4\times4$ matrix $X$ that commutes with all the $\sigma_i\otimes\id_L$. Use the
basis $\ket{00},\ket{01},\ket{10},\ket{11}$, where the first digit belongs to $R$ and the second to $L$. Rather
than solving sixteen equations one by one, write $X$ as four $2\times2$ blocks:
\[
X = \begin{pmatrix} X_{11} & X_{12} \\ X_{21} & X_{22} \end{pmatrix}, \qquad X_{ij} \in M_2(\mathbb C) .
\]
In this basis $\sigma_z\otimes\id_L=\begin{psmallmatrix}\id_2&0\\0&-\id_2\end{psmallmatrix}$ and
$\sigma_x\otimes\id_L=\begin{psmallmatrix}0&\id_2\\\id_2&0\end{psmallmatrix}$. Since $\sigma_x$ and $\sigma_z$
generate all $2\times2$ matrices (their product is $\sigma_x\sigma_z=-i\sigma_y$), it is enough to commute $X$
with these two. First,
\begin{align}
[\sigma_z\otimes\id_L,\, X]
&\eqstep{1} \begin{pmatrix} X_{11} & X_{12} \\ -X_{21} & -X_{22} \end{pmatrix}
- \begin{pmatrix} X_{11} & -X_{12} \\ X_{21} & -X_{22} \end{pmatrix}
\eqstep{2} \begin{pmatrix} 0 & 2X_{12} \\ -2X_{21} & 0 \end{pmatrix} .
\notag
\end{align}
\stepwhy{1}{multiplying by $\sigma_z\otimes\id_L$ on the left flips the sign of the bottom row of blocks;
multiplying on the right flips the sign of the right column of blocks.}\quad
\stepwhy{2}{subtract block by block.}

Setting this to zero forces $X_{12}=X_{21}=0$, so $X$ is block-diagonal. Next, for block-diagonal $X$,
\begin{align}
[\sigma_x\otimes\id_L,\, X]
&\eqstep{1} \begin{pmatrix} 0 & X_{22} \\ X_{11} & 0 \end{pmatrix}
- \begin{pmatrix} 0 & X_{11} \\ X_{22} & 0 \end{pmatrix}
\eqstep{2} \begin{pmatrix} 0 & X_{22}-X_{11} \\ X_{11}-X_{22} & 0 \end{pmatrix} .
\notag
\end{align}
\stepwhy{1}{multiplying by $\sigma_x\otimes\id_L$ on the left swaps the two rows of blocks; multiplying on
the right swaps the two columns of blocks.}\quad
\stepwhy{2}{subtract block by block.}

Setting this to zero forces $X_{22}=X_{11}$. Call this common block $B$. Every $X$ that commutes with $\M$
therefore has the form
\[
X = \begin{pmatrix} B & 0 \\ 0 & B \end{pmatrix} = \id_R \otimes B , \qquad B \in M_2(\mathbb C) .
\]
Since $B$ is an arbitrary $2\times2$ matrix, $\M' = \id_R \otimes B(\HH_L)$. The commutant of ``everything
$R$ can do'' is ``everything $L$ can do''.

Now compute the second commutant $\M''=(\M')'$. An operator
$Y=\begin{psmallmatrix} Y_{11} & Y_{12} \\ Y_{21} & Y_{22} \end{psmallmatrix}$ in $\M''$ must commute with
every $\id_R\otimes B=\begin{psmallmatrix} B & 0 \\ 0 & B \end{psmallmatrix}$. Multiplying out, this says
$Y_{ij}B=BY_{ij}$ for each block and for every $2\times2$ matrix $B$. The only $2\times2$ matrices that commute
with every $2\times2$ matrix are multiples of the identity. (Commuting with
$\begin{psmallmatrix}1&0\\0&0\end{psmallmatrix}$ removes the off-diagonal entries, and commuting with
$\begin{psmallmatrix}0&1\\0&0\end{psmallmatrix}$ makes the two diagonal entries equal.) So
$Y_{ij}=a_{ij}\id_2$, and
\[
Y = \begin{pmatrix} a_{11}\id_2 & a_{12}\id_2 \\ a_{21}\id_2 & a_{22}\id_2 \end{pmatrix}
= \begin{pmatrix} a_{11} & a_{12} \\ a_{21} & a_{22} \end{pmatrix} \otimes \id_L \in \M .
\]
Every such $Y$ lies in $\M$, and every element of $\M$ has this form, so $\M''=\M$. All three algebras are
4-dimensional, and the double commutant theorem holds here by direct calculation.

A second example shows what happens when $\M$ is \emph{not} simply ``all operators on one tensor factor.''
It also previews the notion of a \emph{center}, defined just below. Take the same $\HH=\mathbb C^4$, but now
ignore any tensor-product structure. Let $\M$ be spanned by the two matrices $P_1=\mathrm{diag}(1,1,0,0)$ and
$P_2=\mathrm{diag}(0,0,1,1)$. This is the algebra of an observer who can only tell whether the system is in
``block one'' or ``block two,'' and nothing finer. Solving the same kind of linear system shows that $\M'$ is
8-dimensional. It consists of every operator that acts as an arbitrary $2\times2$ block within each of the
two sectors, without mixing them: two independent $2\times2$ complex matrices, so $4+4=8$ complex
parameters. The second commutant $\M''$ comes out to be exactly $\M$ again, which is 2-dimensional. Something
happens here that did not happen in the first example: $\M\subset\M'$. Both $P_1$ and $P_2$ commute with
\emph{everything} in the block-diagonal algebra $\M'$. So $\M$ is small enough to commute with its own
complement. Whenever this happens, $\M$ has a nontrivial \emph{center}, and that has a physical meaning. $\M$
does not describe one indivisible subsystem. It describes a classical label (which block you are in),
possibly with further subsystems inside each block. The discussion of the center below comes back to this
example.
\end{quote}

\begin{quote}
\textbf{Where the entangled subsystem lives.} Putting the last two subsections together: a subsystem is a von
Neumann algebra $\M\subset B(\HH)$, and its complement is the commutant $\M'$. In this framework, bipartite
entanglement is a statement about the pair $(\M,\M')$, not about a tensor factorization. When a factorization
$\HH=\HH_R\otimes\HH_L$ does exist, the choice $\M=B(\HH_R)\otimes\id_L$ gives back exactly the familiar
story. That is the first worked example above. Chapter~3 calls this special case a \emph{type I factor}.
\end{quote}

Several further facts about commutants are used freely later, so we collect them here. Let $\Alg$ be any
$*$-subalgebra, not necessarily weakly closed.
\begin{itemize}
\item $\Alg'=\Alg'''$. The commutant of $\Alg$ is already a von Neumann algebra: a short argument shows
$\Alg'=(\Alg')''$ directly.
\item $\M\equiv\Alg''$ is the smallest von Neumann algebra containing $\Alg$, and it has the same commutant
as $\Alg$: $\M'=\Alg'''=\Alg'$.
\item For two von Neumann algebras $\M_1,\M_2$, the \textbf{join} $\M_1\vee\M_2\equiv(\M_1\cup\M_2)''$ is the
smallest von Neumann algebra containing both. The \textbf{meet} $\M_1\wedge\M_2\equiv\M_1\cap\M_2$ is their
intersection. The intersection is already a von Neumann algebra, so no double commutant is needed on that
side.
\item These satisfy a De~Morgan-type duality, $(\M_1\vee\M_2)'=\M_1'\wedge\M_2'$. Joining two algebras and
then taking the complement gives the same result as taking each complement first and then intersecting. (The
left side is $(\M_1\cup\M_2)'''=(\M_1\cup\M_2)'$, and an operator commutes with everything in
$\M_1\cup\M_2$ exactly when it commutes with $\M_1$ and with $\M_2$.)
\end{itemize}

The \textbf{center} of $\M$ is $Z(\M)\equiv\M\cap\M'$. It consists of the operators that commute with
everything in $\M$ \emph{and} with everything in $\M'$. In the second worked example above, $\M\subset\M'$,
so $Z(\M)=\M$: the center is as large as it can possibly be. At the other extreme, $\M$ is called a
\textbf{factor} (or is said to be \textbf{primary}) if its center is as small as possible,
$Z(\M)=\mathbb C\,\id$, containing only multiples of the identity. The first worked example,
$\M=B(\HH_R)\otimes\id_L$, is a factor. Indeed, an operator that has both the form $A\otimes\id_L$ and the
form $\id_R\otimes B$ must be a multiple of the identity, as a direct comparison of the two forms shows. A
general von Neumann algebra can always be decomposed into a direct sum, or more generally a direct integral,
of factors. So the classification of von Neumann algebras reduces to classifying factors, plus bookkeeping
for the classical labels carried by the center. For this reason, from here on and for almost all of the
notes, we restrict attention to factors.

\section{Weights and states}

\subsection*{The definitions}

A \textbf{linear functional} $\omega$ on $\M$ is a rule that assigns a complex number $\omega(A)$ to every
$A\in\M$ and respects addition and multiplication by numbers: $\omega(aA+bB)=a\,\omega(A)+b\,\omega(B)$. It is
\textbf{positive} if $\omega(A^\dagger A)\ge0$ for every $A$. This is a natural requirement. The operator
$A^\dagger A$ plays the role of $|A|^2$: it is always a non-negative operator, the operator analogue of a
squared magnitude, so it should never be assigned a negative number. A positive functional is automatically
compatible with the adjoint, $\omega(A^\dagger)=\omega(A)^*$. (This follows from positivity applied to
$(\id+A)^\dagger(\id+A)$ and $(\id+iA)^\dagger(\id+iA)$.) We use this property freely. In these notes a positive
linear functional is called a \textbf{weight}. (In the mathematical literature a weight is also allowed to
take the value $+\infty$ on some positive elements, as the ordinary trace does on $B(\HH)$ when $\HH$ is
infinite-dimensional. Traces of this kind appear in Chapter~3.) A weight is a \textbf{state} if it is also
normalized, $\omega(\id)=1$. It is a \textbf{trace} if $\omega(AB)=\omega(BA)$. This is the cyclic property of
the ordinary matrix trace, $\Tr(AB)=\Tr(BA)$. Positivity and linearity together give a useful inequality for
free. It is the direct analogue of the Cauchy--Schwarz inequality you know from vector inner products:
\begin{equation}
|\omega(A^\dagger B)|^2\le\omega(A^\dagger A)\,\omega(B^\dagger B) .
\label{eq:CS-omega}
\end{equation}

A state $\omega$ is one rule doing one thing: it assigns a single number, $\omega(A)$, to each $A\in\M$. What
that number means depends only on which operator you feed it, not on any change to $\omega$ itself. Feed it
a general observable $A$, and $\omega(A)$ is the expectation value of $A$. Feed it a projection $P_I$
instead, and $\omega(P_I)$ is a probability. Here $P_I$ is the spectral projection onto some range $I$ of
possible outcomes, defined later in this chapter. Positivity and normalization already force
$0\le\omega(P_I)\le1$, because both $P_I=P_I^\dagger P_I$ and $\id-P_I=(\id-P_I)^\dagger(\id-P_I)$ are of the
form $X^\dagger X$. Linearity gives $\omega(P_{I_1})+\omega(P_{I_2})=\omega(P_{I_1\cup I_2})$ for disjoint
ranges. It is the same formula and the same $\omega$ in both cases. The two readings, expectation value and
probability, come entirely from what sits inside the parentheses.

None of these definitions mentions a Hilbert space, a vector, or a density operator. This is the reason for
introducing weights and states this way, rather than going straight to a density matrix. $\omega$ is defined
only as a number assigned to each element of the algebra $\M$. A standard result, stated here without proof,
connects this back to density matrices. Call a state \emph{normal} if it is continuous in a suitable weak
sense; these are the physically sensible states. Every normal state $\omega$ on $\M\subset B(\HH)$ can be
written as $\omega(A)=\Tr(\rho A)$ for some density operator $\rho$ on $\HH$. This $\rho$ acts on all of
$\HH$. It need not belong to $\M$, and in general it is not unique. So the density operator is a consequence
of the definition of $\omega$, not an input to it. $\omega$ is \textbf{faithful} if $\omega(A^\dagger A)=0$
forces $A=0$: nothing nonzero is invisible to it. $\omega$ is \textbf{pure} if it cannot be written as a
nontrivial mixture $\omega=\lambda\omega_1+(1-\lambda)\omega_2$ of two different states, with $0<\lambda<1$. If
it can, $\omega$ is said to \textbf{dominate} $\lambda\omega_1$ and $(1-\lambda)\omega_2$, since
$\omega-\lambda\omega_1$ and $\omega-(1-\lambda)\omega_2$ are still positive.

\subsection*{What happens to the wavefunction once the algebra is taken as primary}

These definitions look innocent. But they demote the wavefunction from a fundamental object to a derived
one, so we unpack them slowly.

An ordinary quantum-mechanical state is a vector $\ket\psi$ or, more generally, a density matrix $\rho$.
Both of these live in, or act on, a Hilbert space. You cannot even write down $\ket\psi$ without first having
$\HH$. The definition of $\omega$ needs none of that. It is a number assigned to each algebra element, and
the algebra $\M$ can be specified abstractly, as a $C^*$-algebra, with no Hilbert space in the definition.
This is more than a cosmetic change. It changes which object is treated as fundamental. In the ordinary
picture, you start with $\HH$, pick a vector $\ket\psi\in\HH$, and compute $\braket{\psi|A|\psi}$ for
whatever operators you like. In the algebraic picture, you start with the algebra $\M$ (what can be measured
or done) and a state $\omega$ on it (an expected outcome for every element of $\M$). Only afterward, if you
want one, do you build a Hilbert space and a vector that reproduce $\omega$. That construction is the GNS
construction, covered in full later in this chapter. The vector it produces is not fundamental data. It is a
derived bookkeeping device, specific to the state $\omega$ you started with. A different state can produce a
different Hilbert space that is physically inequivalent to the first.

The linearity of $\omega$, $\omega(A+B)=\omega(A)+\omega(B)$, also tells you something. It says that $\omega$
describes an \emph{average} over many repetitions, not a single outcome on one system. A single measurement
of $A+B$ on one system does not obviously split into ``the result you would have got for $A$'' plus ``the
result you would have got for $B$.'' In general you cannot even measure $A$ and $B$ together unless they
commute. But the \emph{average} of $A+B$ over a large ensemble of identically prepared systems does split
this way, simply because averaging is linear. So $\omega$ (equivalently the density matrix $\rho$, or in the
pure case the vector $\ket\psi$) is best read as a statement about an ensemble, not about one individual
system.

The statistical algebraic approach of Slavnov~\cite{Slavnov2001} takes this one level deeper. It also
addresses the obvious next question: what, then, describes one individual system? Slavnov starts from
something more primitive than $\omega$. It is a functional $\varphi(\hat A)$ that returns the number a single
measurement of $\hat A$, on this particular system, would read out. It is not defined linearly on the whole
algebra. It is defined on one maximal set of \emph{mutually compatible} observables at a time, that is,
observables that commute and so can be measured jointly. There is no averaging in this definition: $\varphi$
is a record of one concrete measurement outcome. The key point is this. Measuring one observable can disturb
your ability to measure another one that does not commute with it. So no experiment can pin down $\varphi$
completely. An experiment can only tell you that $\varphi$ belongs to an \emph{equivalence class} of
functionals that all agree with your data on the compatible set of observables $\{Q\}$ you chose to measure.
Slavnov proposes that this equivalence class, and not $\varphi$ itself, is what a conventional ``quantum
state'' $\Psi_Q$ is. On this view a quantum state is not a container of hidden, definite properties. It is a
precise record of how much you could possibly know, given that measuring one thing costs you the ability to
measure something incompatible with it.

On this view, one idea accounts for several facts that otherwise look like separate puzzles:

\begin{itemize}
\item \textbf{Why quantum probability is not ``we just don't know the hidden variable yet.''} The same
individual $\varphi$ can belong to the class $\Psi_Q=\{\varphi\}_Q$ or to a different class
$\Psi_P=\{\varphi\}_P$. Which one depends on which compatible set, $\{Q\}$ or $\{P\}$, the experimenter chooses
to measure. So the ``quantum state'' assigned to the system depends on an experimental choice made after the
system was prepared. Slavnov identifies this as the root of the Einstein--Podolsky--Rosen puzzle. On his
account, the EPR puzzle is not an additional strange feature of quantum mechanics. It is the same fact seen
from another angle.
\item \textbf{Why $\rho$, not $\ket\psi$, is what enters every physical prediction.} $\rho$ (equivalently
$\omega$) is, by construction, one step of ensemble averaging removed from an individual $\varphi$. Textbooks
such as Sakurai's replace $\ket\psi$ by $\rho=\ket\psi\bra\psi$ as the basic object, so that pure and mixed
states are handled in the same way. That is the first half of this demotion. The algebraic framework of these
notes completes it: it replaces $\rho$, which needs a Hilbert space to be written as a matrix, with $\omega$,
which does not.
\item \textbf{Why a Hilbert space appears at all, and in what sense it comes second.} Slavnov states this
directly: \emph{Hilbert space vectors (wavefunctions), the operators in the Hilbert space, and other purely
mathematical concepts are not primary elements. These concepts only appear at the second stage of the
formulation.} The primary elements are the algebra of observables and the state, both tied directly to
experiment. The state can be taken in Slavnov's individual sense $\varphi$, or in the ensemble sense $\omega$
used throughout these notes. The Hilbert space, its vectors, and the operators acting on them are all built
afterward, mechanically, by the GNS construction applied to $\omega$.
\end{itemize}

To state the conclusion of this viewpoint plainly, in one place: the wavefunction is not a physical entity
that stores values the way a hard drive stores bits. It is the name for an equivalence class of individual
measurement records. The records in one class cannot be told apart, because no experiment could distinguish
them, given what was actually chosen to be measured. The algebra, which says what is measurable and how
measurements combine, is the primary, observer-independent structure. The wavefunction depends on a choice
of what you decided to look at. This is an interpretation, not a theorem, and the mathematics in the rest of
the notes do not depend on adopting it. What the rest of the notes do use is the weaker statement: the
algebra and the state come first, and the Hilbert space is built from them.

\section{\texorpdfstring{$C^*$}{C*} and von Neumann algebras generated by a single operator}

This section is short. Given the two definitions ($C^*$-algebra and von Neumann algebra) from earlier in this
chapter, it is self-contained. It gives the concrete reason why projections, the subject of the next
section, live in von Neumann algebras and generally not in the smaller $C^*$-algebras.

Take one bounded, self-adjoint operator $\mathcal O$. What is the smallest $C^*$-algebra containing it,
$\Alg(\mathcal O)$, and the smallest von Neumann algebra containing it, $\M(\mathcal O)$? Both contain every
polynomial $\sum_n c_n\mathcal O^n$ in $\mathcal O$, since polynomials are made by repeated multiplication and
addition. Each must then be completed, in its own sense of limit, to become a full $C^*$-algebra or von
Neumann algebra.

The \textbf{spectrum} $\sigma(\mathcal O)$ of $\mathcal O$ is the set of numbers $\kappa$ for which
$\mathcal O-\kappa\id$ is not invertible. For a finite Hermitian matrix, the spectrum is its set of
eigenvalues: the numbers a measurement of $\mathcal O$ can return. A standard fact of functional analysis
(stated here without proof) is that $\Alg(\mathcal O)$ is isomorphic to the algebra of \emph{continuous}
functions on $\sigma(\mathcal O)$. Every element of $\Alg(\mathcal O)$ can be written as $f(\mathcal O)$ for a
continuous function $f$, and every such $f(\mathcal O)$ lies in $\Alg(\mathcal O)$. A simple example makes
this plausible. Let $\mathcal O$ be a finite Hermitian matrix with distinct eigenvalues
$\kappa_1,\dots,\kappa_n$. Applying a polynomial, or any function, to $\mathcal O$ means applying it to each
eigenvalue and leaving the eigenvectors alone. So specifying $f(\mathcal O)$ is the same as specifying the $n$
numbers $f(\kappa_1),\dots,f(\kappa_n)$, that is, a function on the finite set $\{\kappa_1,\dots,\kappa_n\}$.
On a finite set every function is continuous, so this example cannot yet show the distinction that matters
below. It does show the general shape of the isomorphism.

\begin{physconnect}[: Eigenvalues vs.\ Eigenvectors]
The definition of $\sigma(\mathcal O)$ never mentions an eigenvector or a Hilbert space. A number $\lambda$
lies in $\sigma(\mathcal O)$ because $\mathcal O-\lambda\id$ has no inverse inside the algebra. So eigenvalues
do not need to be demoted, as the vacuum vector $\ket\Omega$ will be in the GNS construction later in this
chapter. They are already an algebra-level notion, fixed before any state or representation is chosen. What
is demoted is the eigen\emph{vector}. An eigenvector lives in whatever Hilbert space GNS builds for a given
$\omega$, so it appears only once $\omega$ is chosen.

Check this on the number operator $N=a^\dagger a$ of the harmonic oscillator. (Strictly, $N$ is unbounded, so
it is not an element of a $C^*$-algebra. The argument below uses only the commutation relation and
positivity, so this does not matter.) We show that the possible eigenvalues of $N$ are fixed by
$[a,a^\dagger]=1$ together with positivity, in every representation at once. Suppose $Nv=\lambda v$ for some
nonzero vector $v$, in any representation of the algebra on a Hilbert space. First, $\lambda\ge0$:
\begin{align}
\lambda\,\|v\|^2
&\eqstep{1} \braket{v|Nv} \notag\\
&\eqstep{2} \braket{av|av} = \|av\|^2 \ \ge\ 0 . \notag
\end{align}
\stepwhy{1}{$Nv=\lambda v$ by hypothesis.}\quad
\stepwhy{2}{$N=a^\dagger a$ and the definition of the adjoint.}

Next, lower with $a$:
\begin{align}
N(av)
&\eqstep{1} \big(aN+[N,a]\big)v \notag\\
&\eqstep{2} a(\lambda v) + (-a)v \notag\\
&\eqstep{3} (\lambda-1)(av) . \notag
\end{align}
\stepwhy{1}{$Na=aN+[N,a]$, the definition of the commutator.}\quad
\stepwhy{2}{$Nv=\lambda v$ by hypothesis, and $[N,a]=[a^\dagger a,a]=[a^\dagger,a]\,a=-a$ from $[a,a^\dagger]=1$.}\quad
\stepwhy{3}{collect the two terms, both proportional to $av$.}

So $av$ is either zero or an eigenvector with eigenvalue $\lambda-1$. Repeating, $a^kv$ is either zero or an
eigenvector with eigenvalue $\lambda-k$. Every eigenvalue is $\ge0$, so the vectors $a^kv$ cannot all be
nonzero. Let $k$ be the last index with $a^kv\ne0$. Then $a(a^kv)=0$, so $N(a^kv)=a^\dagger a(a^kv)=0$, and
the eigenvalue $\lambda-k$ must be $0$. Hence $\lambda=k\in\{0,1,2,\dots\}$. The same computation with
$[N,a^\dagger]=a^\dagger$ gives $N(a^\dagger v)=(\lambda+1)(a^\dagger v)$. Here $a^\dagger v$ is never zero:
$a^\dagger v=0$ would give $aa^\dagger v=(N+1)v=0$, so $\lambda=-1$, which contradicts $\lambda\ge0$. So once
one eigenvalue occurs, all of $0,1,2,\dots$ occur. These are the quantized levels of the oscillator. They
follow from the multiplication table of the algebra and positivity, whichever Hilbert space the operators
act on. Choosing a state $\omega$ and running GNS only decides which eigenvectors appear, and with what
weights. It never changes the spectrum itself.
\end{physconnect}

$\M(\mathcal O)$, by contrast, is isomorphic to an algebra of \emph{bounded} functions on $\sigma(\mathcal O)$
that need not be continuous. (Precisely: bounded measurable functions, where two functions count as equal if
they differ only on a set that the spectral measure of $\mathcal O$ does not see.) Every continuous function
on the compact set $\sigma(\mathcal O)$ is bounded, so $\Alg(\mathcal O)\subset\M(\mathcal O)$. This fits with
von Neumann algebras being the larger, more complete objects. The difference becomes concrete and important
when $\sigma(\mathcal O)$ is a continuum, such as an interval of real numbers, rather than a finite set of
isolated points. This happens when $\mathcal O$ has continuous spectrum, which is the generic case in field
theory and in the $N\to\infty$ limits these notes are about. On a continuum, the \textbf{indicator function} of a
subset $I\subset\sigma(\mathcal O)$,
\[
f_I(\kappa) = \begin{cases} 1 & \kappa\in I \\ 0 & \kappa\notin I \end{cases} ,
\]
is bounded, since it only takes the values $0$ and $1$. But it is discontinuous wherever $I$ has an edge
inside the continuum. The corresponding operator $P_I\equiv f_I(\mathcal O)$ is a \textbf{spectral
projection}. For a matrix, it projects onto the span of the eigenvectors whose eigenvalues lie in $I$. In
general, it is the operator that answers the yes/no question ``does a measurement of $\mathcal O$ give a value
in the range $I$?'' Indicator functions are generically discontinuous on a continuum, so spectral projections
generically lie in $\M(\mathcal O)$ but not in $\Alg(\mathcal O)$. This is the concrete, checkable reason why
von Neumann algebras, and not merely $C^*$-algebras, are the right setting for the classification in the next
two sections. That classification is built entirely out of projections, and a $C^*$-algebra generally does
not contain the sharp yes/no measurement operators it needs.

\section{Projections}

\subsection*{Definitions, and what they mean physically}

A \textbf{projection} is a self-adjoint operator $P$ with $P^2=P$. Applying it twice does nothing more than
applying it once. This is the algebraic signature of ``select a subspace and discard everything orthogonal to
it,'' which is what a sharp yes/no measurement does. The spectral theorem writes any self-adjoint operator in
terms of its spectral projections, as discussed in the previous section. As a result, every element of a von
Neumann algebra is a norm limit of finite linear combinations of projections in the algebra. This is why the
classification that follows can be phrased entirely in terms of projections, with nothing lost.

\begin{physconnect}[: Probabilities from Projections]
The section on weights and states noted that $\omega(P_I)$ reads as a probability because $P_I$ is a
projection. Here is where that fact comes from. The previous section fixed the possible outcomes of measuring
$\mathcal O$ as its spectrum $\sigma(\mathcal O)$. This is an algebra-level fact, decided before any state is
chosen. For a range $I\subset\sigma(\mathcal O)$, the spectral projection $P_I$, built from $\mathcal O$ by
the indicator-function construction, is itself an element of the von Neumann algebra. The Born rule of
ordinary quantum mechanics gives the probability of an outcome in $I$ as $\braket{\psi|P_I|\psi}$. The
algebraic version replaces $\ket\psi$ by $\omega$, so the probability is $\omega(P_I)$. The probability of an
outcome is therefore not a second thing that $\omega$ supplies on top of expectation values. It is $\omega$
evaluated on one particular algebra element, the spectral projection for that outcome, by the same rule
$\omega$ uses for everything else.
\end{physconnect}

\begin{quote}
\textit{Worked example: spectral projections and the Born rule on a single qubit.} Let
$A = \sigma_x + \sigma_z = \begin{psmallmatrix} 1 & 1 \\ 1 & -1 \end{psmallmatrix}$, an observable on
$\mathbb C^2$. We construct its spectral projections explicitly and use them to compute probabilities.
\begin{enumerate}
\item \textbf{Eigenvalues.} The characteristic equation is $\det(A - \lambda\id) = \lambda^2 - 2 = 0$, so
the eigenvalues are $\lambda_\pm = \pm\sqrt{2}$.
\item \textbf{Spectral projections by Lagrange interpolation.} Let a $2\times2$ Hermitian matrix $A$ have
distinct eigenvalues $\lambda_1, \lambda_2$. The projection onto the eigenspace of $\lambda_1$ is
$P_1 = (A - \lambda_2\id)/(\lambda_1 - \lambda_2)$. The reason is that the polynomial
$f(x)=(x-\lambda_2)/(\lambda_1-\lambda_2)$ equals $1$ at $\lambda_1$ and $0$ at $\lambda_2$, so $f(A)$ is the
indicator function of $\{\lambda_1\}$ applied to $A$. For our $A$ this gives
\begin{align}
P_+ &= \frac{A + \sqrt2\,\id}{\sqrt2 + \sqrt2}
= \frac{1}{2\sqrt2}\begin{pmatrix} 1+\sqrt2 & 1 \\ 1 & \sqrt2-1 \end{pmatrix} , \notag\\
P_- &= \frac{A - \sqrt2\,\id}{-\sqrt2 - \sqrt2}
= \frac{1}{2\sqrt2}\begin{pmatrix} \sqrt2-1 & -1 \\ -1 & \sqrt2+1 \end{pmatrix} . \notag
\end{align}
\item \textbf{Checking the projection properties.}
\begin{itemize}
\item $P_+ + P_- = \frac{1}{2\sqrt2}\begin{psmallmatrix} 2\sqrt2 & 0 \\ 0 & 2\sqrt2 \end{psmallmatrix} = \id_2$.
\item $P_+ P_- = \frac{1}{8}\begin{psmallmatrix} 1+\sqrt2 & 1 \\ 1 & \sqrt2-1 \end{psmallmatrix}
\begin{psmallmatrix} \sqrt2-1 & -1 \\ -1 & \sqrt2+1 \end{psmallmatrix} = 0$. For example, the top-left entry
is $\tfrac18\big[(1+\sqrt2)(\sqrt2-1)-1\big]=\tfrac18(1-1)=0$.
\item $P_+^2 = P_+$ and $P_-^2 = P_-$. This follows from the two previous lines:
$P_+^2=P_+(\id-P_-)=P_+-P_+P_-=P_+$, and in the same way for $P_-$.
\item Spectral decomposition: $\lambda_+ P_+ + \lambda_- P_- = \sqrt2\, P_+ - \sqrt2\, P_-
= \begin{psmallmatrix} 1 & 1 \\ 1 & -1 \end{psmallmatrix} = A$.
\end{itemize}
\item \textbf{Born-rule probabilities.} Suppose the system is in the state
$\ket\psi = \ket0 = \begin{psmallmatrix} 1 \\ 0 \end{psmallmatrix}$. The probabilities of the two outcomes are
\begin{align}
\operatorname{Prob}(+\sqrt2) &= \braket{0|P_+|0} = \frac{1+\sqrt2}{2\sqrt2}
= \frac12 + \frac{1}{2\sqrt2} \approx 0.853553 , \notag\\
\operatorname{Prob}(-\sqrt2) &= \braket{0|P_-|0} = \frac{\sqrt2-1}{2\sqrt2}
= \frac12 - \frac{1}{2\sqrt2} \approx 0.146447 . \notag
\end{align}
They add up to exactly $1$, as they must, since $P_++P_-=\id$. The expectation value is
\begin{align}
\braket{A}
&\eqstep{1} (+\sqrt2)\operatorname{Prob}(+\sqrt2) + (-\sqrt2)\operatorname{Prob}(-\sqrt2) \notag\\
&\eqstep{2} \sqrt2\Big(\frac12 + \frac{1}{2\sqrt2}\Big) - \sqrt2\Big(\frac12 - \frac{1}{2\sqrt2}\Big) \notag\\
&\eqstep{3} 2\cdot\frac{\sqrt2}{2\sqrt2} = 1 \notag\\
&\eqstep{4} \braket{0|\sigma_x+\sigma_z|0} . \notag
\end{align}
\stepwhy{1}{take $\braket{0|\cdot|0}$ of the spectral decomposition $A=\sqrt2\,P_+-\sqrt2\,P_-$.}\quad
\stepwhy{2}{insert the two probabilities just computed.}\quad
\stepwhy{3}{the two $\sqrt2\cdot\tfrac12$ terms cancel, and the two remaining terms are equal.}\quad
\stepwhy{4}{$\braket{0|\sigma_x|0}=0$ and $\braket{0|\sigma_z|0}=1$.}

Numerically, $\sqrt2\,(0.853553-0.146447)=\sqrt2\,(0.707107)=1.000000$.
\end{enumerate}
The example shows how spectral projections break an observable into exact, mutually orthogonal yes/no
questions. Each question has a probability, here $\braket{0|P|0}$, and the expectation value is recovered by
weighting each outcome by its probability.
\end{quote}

There is an exact correspondence between projections in $B(\HH)$ and closed subspaces of $\HH$. A projection
$P$ picks out the subspace $P\HH$, which consists of everything you get by applying $P$ to vectors in $\HH$.
Conversely, every closed subspace has a unique projection onto it. If $P\in\M$, the subspace $P\HH$ is said to
``belong to $\M$.'' Physically, it is a subspace that an observer with access to $\M$ can single out by a
measurement. The largest projection is the identity $\id$ (the whole space). Projections are partially
ordered: $P\le Q$ means $PQ=P$, or equivalently $P\HH\subseteq Q\HH$, so the subspace of $P$ sits entirely
inside the subspace of $Q$.

Two projections $P,Q\in\M$ are called \textbf{Murray--von Neumann equivalent}, written $P\sim Q$, if there is
a \textbf{partial isometry} $V\in\M$ with $P=V^\dagger V$ and $Q=VV^\dagger$. A partial isometry is an operator
that maps its input subspace onto its output subspace without stretching or shrinking lengths. So the
condition says that $V$ maps the subspace $P\HH$ isometrically (preserving lengths) onto the subspace
$Q\HH$, and that $V$ itself is available inside $\M$. This is the algebra-relative version of ``these two
subspaces have the same size.'' It is not the ordinary linear-algebra notion of dimension, which looks only
at the subspace itself. It is tied to what an observer with access to $\M$ can actually do to compare one
subspace with another.

\begin{quote}
\textit{Worked example.} Return to $\M=B(\HH_R)\otimes\id_L$ on $\HH=\mathbb C^2\otimes\mathbb C^2$ from the
double-commutant example above. Let $P=\ket0_R\!\bra0_R\otimes\id_L$ and $Q=\ket1_R\!\bra1_R\otimes\id_L$.
Each is a rank-2 projection on the full 4-dimensional $\HH$. It picks out a 2-dimensional subspace: the
$R$-qubit is fixed to $\ket0$ or to $\ket1$, and the $L$-qubit can be anything. The operator
$V=\ket1_R\!\bra0_R\otimes\id_L$ is in $\M$, since it has the form $A\otimes\id_L$ with $A=\ket1\bra0$.
Direct matrix multiplication gives $V^\dagger V=\ket0_R\!\bra0_R\otimes\id_L=P$ and
$VV^\dagger=\ket1_R\!\bra1_R\otimes\id_L=Q$. So $P\sim Q$, through a partial isometry that lies entirely
inside $\M$. This captures the intuitive fact that ``the $R$-qubit is $\ket0$'' and ``the $R$-qubit is
$\ket1$'' are subspaces of the same size, related by a flip that an $R$-observer can actually perform.
\end{quote}

A nonzero projection $P\in\M$ is called \textbf{finite} if it is \emph{not} equivalent to any strictly smaller
projection $Q<P$ in $\M$. It is called \textbf{infinite} if it is equivalent to some strictly smaller
projection $Q<P$ in $\M$. This definition takes some getting used to, because it does \emph{not} match the
linear-algebra notion of ``finite-dimensional.'' The following example shows the difference. The useful
intuition turns out to be that the smallest projections pick out one copy of an irreducible representation.

Take $\M=B(\HH_1)\otimes\id_2$ on $\HH=\HH_1\otimes\HH_2$, where $\HH_2$ is \emph{infinite}-dimensional. Let
$P=\ket\psi\bra\psi\otimes\id_2$ for a single unit vector $\ket\psi\in\HH_1$. As a subspace of the full
Hilbert space, $P\HH$ is infinite-dimensional: it is a whole copy of $\HH_2$. Yet $P$ is a \emph{finite}
projection of the algebra $\M$. There is no strictly smaller projection in $\M$ equivalent to it, because
$\M$ contains no nonzero projection strictly below $P$ at all. Inside $\M\cong B(\HH_1)$, the projection $P$
corresponds to a rank-one projection on $\HH_1$, the smallest nonzero rank there is. Here is the precise form
of the irreducible-representation intuition. As an \emph{abstract algebra}, forgetting the Hilbert space it
acts on, $\M$ is isomorphic to $B(\HH_1)$. That is a single copy of the algebra of all operators on $\HH_1$,
and it acts irreducibly on $\HH_1$. The factor $\HH_2$ in $\HH=\HH_1\otimes\HH_2$ only counts \emph{how many
copies} of this irreducible representation sit side by side in the big Hilbert space: one copy for each basis
vector of $\HH_2$. The projection $P=\ket\psi\bra\psi\otimes\id_2$ picks out one direction \emph{within} the
irreducible representation, and it does so in every one of the infinitely many copies at once. That is why
$P$ is as small as $\M$ can make anything (finite, and in fact minimal, in the algebra) while being
infinite-dimensional as a subspace of $\HH$. Finiteness in this algebra-relative sense is a statement about
the irreducible representation, not about the dimension of the ambient Hilbert space. Keeping this
distinction in mind is what makes the type classification in the next section make sense, rather than look
like a strange use of the word ``finite.''

A projection $P$ is called \textbf{minimal} in $\M$ if it is nonzero and $\M$ contains no nonzero projection
strictly smaller than $P$. This is the algebra-relative notion of ``as small as a measurement outcome can
be.'' Every minimal projection is finite, since there is nothing smaller for it to be equivalent to. The
converse fails: a projection can be finite without being minimal. For example, in $\M=B(\mathbb C^3)$ every
projection is finite, because equivalent projections have equal rank. But $\mathrm{diag}(1,1,0)$ is not
minimal, since $\mathrm{diag}(1,0,0)$ lies strictly below it. More strikingly, an algebra can have no minimal
projections at all while still having plenty of finite ones. This possibility is the heart of the
classification in the next section.

Now apply these definitions to the largest projection, $P=\id$. For a von Neumann \emph{factor} $\M$, the
identity is either finite or infinite, and this one question sorts factors into two coarse classes. If $\id$
is finite, then every projection in $\M$ is also finite (this can be shown from the definitions; we omit the
proof), and $\M$ is called a \textbf{finite von Neumann factor}. If $\id$ is infinite, $\M$ is called an
\textbf{infinite von Neumann factor}. In that case, when $\HH$ is separable, any two infinite projections in
$\M$ are equivalent to each other, so in particular every infinite projection is equivalent to the identity
itself (also stated without proof).

\section{Classification of von Neumann factors}

\subsection*{Building a dimension function out of nothing but $\sim$ and $\le$}

Here is the payoff of the last two sections. Using only the equivalence relation $P\sim Q$ and the ordering
$P\le Q$, one can prove the following theorem (a standard result, stated here without proof). Every von
Neumann factor $\M$ has a \textbf{dimension function} $d$. It assigns a number $d(P)\in[0,\infty]$ to each
projection $P\in\M$, and it has three properties:
\begin{itemize}
\item $d(P)=d(Q)$ if and only if $P\sim Q$;
\item $d(P+Q)=d(P)+d(Q)$ whenever $P$ and $Q$ are orthogonal, $PQ=0$;
\item $d(P)<\infty$ if and only if $P$ is finite.
\end{itemize}
These properties fix $d$ up to an overall constant that you are free to rescale. In words: equivalent
subspaces get the same size, sizes add, and finite projections get finite sizes. Additivity also makes $d$
respect the ordering. If $P\le Q$, then $Q-P$ is a projection orthogonal to $P$, so
$d(Q)=d(P)+d(Q-P)\ge d(P)$. The inequality is strict when $P<Q$ and $P$ is finite, because $d(Q-P)>0$ for the
nonzero projection $Q-P$. These are exactly the properties you would want from anything worth calling a
``size.'' The function $d$ can also be packaged as a \textbf{trace}, $\tr(P)\equiv d(P)$, extended from
projections to other elements of $\M$ by linearity, using the fact that every self-adjoint operator is built
from its spectral projections. When $d(\id)=\infty$, this $\tr$ is finite only on part of $\M$, just as the
ordinary trace on an infinite-dimensional Hilbert space is finite only for trace-class operators. This $\tr$
plays the role of the ordinary matrix trace, generalized to an abstract algebra with no particular Hilbert
space singled out.

Why is the ordinary formula ``dimension $=\Tr P$'' not the right notion here? Here $\Tr$ is the ordinary
Hilbert-space trace, the sum of diagonal entries. The answer shows why an algebra-relative $d(P)$ is needed
at all. Take again $P=\ket\psi\bra\psi\otimes\id_2$ on $\HH_1\otimes\HH_2$ with $\HH_2$ infinite-dimensional,
the same example as above. The Hilbert-space trace $\Tr P$ is infinite, because $P\HH$ is an
infinite-dimensional subspace of $\HH$. But from the point of view of the algebra $\M\cong B(\HH_1)$, the
projection $P$ is the smallest nonzero thing there is: rank one, within $\HH_1$. The ordinary trace $\Tr$
sees the raw size of the ambient Hilbert space, and for describing what $\M$ can distinguish it gives the
wrong answer. The algebra-relative $d(P)$, normalized so that $d(P)=1$ for this minimal $P$, correctly reports
``this is the smallest measurable unit $\M$ has access to,'' however large the corresponding subspace of $\HH$
is. This is the concrete sense in which $\tr$ is a \emph{renormalized} trace. It strips away the part of $\HH$
that $\M$ never had access to, and it reports only the size that $\M$ itself can distinguish.

\subsection*{The three types}

The whole classification now follows from one question: what values can $d(P)$ take, as $P$ ranges over all
the projections in $\M$?

\begin{itemize}
\item \textbf{Type I: $\M$ has minimal projections.} Normalize $d$ so that a minimal projection has $d=1$.
You are always free to rescale $d$ by an overall constant, so this is only a choice of units, like choosing
what counts as one unit of length. Every projection is then a sum of mutually orthogonal minimal
projections, and it can be shown that $d(P)$ is a nonnegative integer or $\infty$ for every $P$. This is the
familiar notion of dimension, now derived from the algebra instead of assumed. If $n=d(\id)$ is finite, $\M$
is called type $\mathrm I_n$. If $d(\id)=\infty$, so that there is no upper bound on the number of mutually
orthogonal minimal projections, $\M$ is type $\mathrm I_\infty$. For an ordinary separable Hilbert space
$\HH$, the algebra $B(\HH)$ itself is type $\mathrm I_n$ if $\HH$ is $n$-dimensional, and type
$\mathrm I_\infty$ if $\HH$ is infinite-dimensional with a countable basis. In this case $\tr$ is the ordinary
matrix trace $\Tr_\HH$, with no renormalization needed, because there is no part of $\HH$ that the algebra
fails to see.

The most important fact tying this back to ordinary quantum mechanics is the following. $\M$ is a type I
factor \emph{if and only if} there is a Hilbert space factorization $\HH=\HH_R\otimes\HH_L$ with
\[
\M=B(\HH_R)\otimes\id_L, \qquad \tr=\Tr_{\HH_R}, \qquad \M'=\id_R\otimes B(\HH_L) .
\]
This result closes the type classification, and Chapter~3 uses it constantly. Many books state it without
proof. The proof is short enough to give here, and it shows exactly how minimal projections build a tensor
product.

\begin{keyresult}[: Derivation of the Type I Factorization Theorem]
\textbf{Theorem:} A von Neumann algebra $\M \subseteq B(\HH)$ is a type I factor if and only if there exists a
unitary isomorphism $U: \HH \xrightarrow{\sim} \HH_R \otimes \HH_L$ such that
\[
U \M U^\dagger = B(\HH_R) \otimes \id_L, \qquad U \M' U^\dagger = \id_R \otimes B(\HH_L) .
\]
\textbf{Proof ($\implies$):}
\begin{enumerate}[leftmargin=*]
\item \textbf{Minimal projection and orthogonal resolution.}
By the definition of type I, $\M$ contains a minimal projection $P \in \M$. Because $P$ is minimal, the
compressed algebra $P\M P$ contains only multiples of $P$:
\[
P \M P = \mathbb{C} P .
\]
(The reason: $P\M P$ is a von Neumann algebra on $P\HH$ whose only projections are $0$ and $P$, and a von
Neumann algebra is generated by its projections.) Since $\M$ is a factor, its center is trivial,
$\mathcal{Z}(\M) \equiv \M \cap \M' = \mathbb{C}\id$. So the central support of $P$ (the smallest central
projection above $P$) is $c(P) = \id$. By the comparison theorem for projections in a factor, any two minimal
projections are Murray--von Neumann equivalent. By Zorn's lemma, we can choose a maximal family
$\{P_i\}_{i \in I}$ of mutually orthogonal minimal projections, each equivalent to $P$. Maximality together
with $c(P)=\id$ implies that their sum is the identity on $\HH$:
\[
\sum_{i \in I} P_i = \id_\HH, \qquad P_i P_j = \delta_{ij} P_i .
\]

\item \textbf{Equivalence through partial isometries.}
Fix a base index $0 \in I$ with $P_0 \equiv P$. Since $P_i \sim P$, there are partial isometries $V_i \in \M$
with
\[
V_i^\dagger V_i = P, \qquad V_i V_i^\dagger = P_i \qquad (\text{with } V_0 \equiv P).
\]
So $V_i$ maps the base subspace $P\HH$ isometrically onto the orthogonal subspace $P_i\HH$.

\item \textbf{Construction of the unitary $U$.}
Define two Hilbert spaces:
\[
\HH_R \equiv \ell^2(I) \quad \text{with orthonormal basis } \{\ket{i}\}_{i \in I}, \qquad \HH_L \equiv P\HH .
\]
Define the linear map $U: \HH \to \HH_R \otimes \HH_L$ by its action on any vector $\ket\psi \in \HH$:
\[
U \ket\psi \equiv \sum_{i \in I} \ket{i} \otimes \big(V_i^\dagger \ket\psi\big) .
\]
Each $V_i^\dagger \ket\psi = P V_i^\dagger \ket\psi$ lies in $P\HH = \HH_L$, so this is well defined. We check
that $U$ is an isometry:
\begin{align}
\|U\ket\psi\|^2 = \sum_{i \in I} \|V_i^\dagger \ket\psi\|^2
&\eqstep{1} \sum_{i \in I} \braket{\psi | V_i V_i^\dagger | \psi} \notag\\
&\eqstep{2} \sum_{i \in I} \braket{\psi | P_i | \psi}
\ \eqstep{3}\ \Braket{\psi \Big| \sum_{i \in I} P_i \Big| \psi}
\ \eqstep{4}\ \braket{\psi|\psi} . \notag
\end{align}
\stepwhy{1}{$\|v\|^2=\braket{v|v}$ applied to $v=V_i^\dagger\ket\psi$, using $(V_i^\dagger)^\dagger=V_i$.}\quad
\stepwhy{2}{the partial isometry relation $V_iV_i^\dagger=P_i$ from Step~2 above.}\quad
\stepwhy{3}{linearity of the inner product, pulling the sum through $\braket{\psi|\cdot|\psi}$.}\quad
\stepwhy{4}{the resolution of the identity, $\sum_i P_i=\id_\HH$, established in Step~1.}

The adjoint map $U^\dagger: \HH_R \otimes \HH_L \to \HH$ acts on elementary basis tensors as
\[
U^\dagger \big(\ket{i} \otimes \ket{\phi_L}\big) = V_i \ket{\phi_L}, \qquad \ket{\phi_L} \in P\HH .
\]
Computing $U U^\dagger$ on basis vectors:
\begin{align}
U U^\dagger \big(\ket{j} \otimes \ket{\phi_L}\big)
&\eqstep{1} U \big(V_j \ket{\phi_L}\big) \notag\\
&\eqstep{2} \sum_{i \in I} \ket{i} \otimes \big(V_i^\dagger V_j \ket{\phi_L}\big)
\ \eqstep{3}\ \ket{j} \otimes \big(P \ket{\phi_L}\big) = \ket{j} \otimes \ket{\phi_L} . \notag
\end{align}
\stepwhy{1}{the definition of $U^\dagger$ on the basis tensor $\ket j\otimes\ket{\phi_L}$.}\quad
\stepwhy{2}{the definition of $U$ applied to the vector $V_j\ket{\phi_L}$.}\quad
\stepwhy{3}{$V_i^\dagger V_j = V_i^\dagger P_i P_j V_j = \delta_{ij} P$, so only the $i=j$ term in the sum
survives, and $P\ket{\phi_L}=\ket{\phi_L}$ since $\ket{\phi_L}\in P\HH$.}

So $U$ is an isometry with $UU^\dagger=\id$, which means $U$ is unitary.

\item \textbf{Action on the algebra $\M$.}
Let $A \in \M$ be any element. Compute $U A U^\dagger$ on $\ket{j} \otimes \ket{\phi_L}$:
\begin{align}
U A U^\dagger \big(\ket{j} \otimes \ket{\phi_L}\big)
&\eqstep{1} U \big(A V_j \ket{\phi_L}\big) \notag\\
&\eqstep{2} \sum_{i \in I} \ket{i} \otimes \big(V_i^\dagger A V_j \ket{\phi_L}\big) . \notag
\end{align}
\stepwhy{1}{the definition of $U^\dagger$ on basis tensors, then applying $A$.}\quad
\stepwhy{2}{the definition of $U$ applied to the vector $AV_j\ket{\phi_L}$.}

The key object is the operator $V_i^\dagger A V_j$. It lies in $P\M P$:
\[
V_i^\dagger A V_j = (P V_i^\dagger) A (V_j P) = P \big(V_i^\dagger A V_j\big) P \in P \M P .
\]
Because $P$ is minimal, $P \M P = \mathbb{C} P$. Therefore $V_i^\dagger A V_j$ is a multiple of $P$:
\[
V_i^\dagger A V_j = a_{ij} P \quad \text{for some number } a_{ij} \in \mathbb{C} .
\]
Acting on $\ket{\phi_L} \in P\HH$, this gives $V_i^\dagger A V_j \ket{\phi_L} = a_{ij} \ket{\phi_L}$. Thus
\begin{align}
U A U^\dagger \big(\ket{j} \otimes \ket{\phi_L}\big)
&\eqstep{1} \sum_{i \in I} a_{ij} \ket{i} \otimes \ket{\phi_L} \notag\\
&\eqstep{2} \Big(\sum_{i,k \in I} a_{ik} \ket{i}\bra{k} \otimes \id_L \Big) \big(\ket{j} \otimes \ket{\phi_L}\big) . \notag
\end{align}
\stepwhy{1}{$V_i^\dagger AV_j\ket{\phi_L}=a_{ij}\ket{\phi_L}$, just shown.}\quad
\stepwhy{2}{$\braket{k|j}=\delta_{kj}$ picks out exactly the $k=j$ column of the matrix $(a_{ik})$.}

So every operator in $\M$ acts trivially on $\HH_L$ and as a matrix on $\HH_R$. Conversely, every rank-one
operator $\ket{i}\bra{j}\otimes\id_L$ comes from an element of $\M$: a computation like the one above gives
$U V_i V_j^\dagger U^\dagger = \ket{i}\bra{j}\otimes\id_L$, and $V_iV_j^\dagger\in\M$. Linear combinations of
these operators are weakly dense in $B(\HH_R)\otimes\id_L$, and $\M$ is weakly closed. Thus
$U \M U^\dagger = B(\HH_R) \otimes \id_L$.

\item \textbf{The commutant and the trace.}
Conjugating by the unitary $U$ and using the previous step:
\begin{align}
U \M' U^\dagger
&\eqstep{1} (U \M U^\dagger)' \notag\\
&\eqstep{2} \big(B(\HH_R) \otimes \id_L\big)'
\ \eqstep{3}\ \id_R \otimes B(\HH_L) . \notag
\end{align}
\stepwhy{1}{$B$ commutes with every element of $\M$ exactly when $UBU^\dagger$ commutes with every element of
$U\M U^\dagger$, since $U^\dagger U=\id$.}\quad
\stepwhy{2}{$U\M U^\dagger=B(\HH_R)\otimes\id_L$ from Step~4.}\quad
\stepwhy{3}{the commutant of ``all operators on one tensor factor'' is ``all operators on the other,'' as
in the first double-commutant example earlier in this chapter.}

The trace $\tr$ on $\M$, normalized so that minimal projections have trace $1$, corresponds under $U$ to the
standard trace $\Tr_{\HH_R}$ on $B(\HH_R)$.

\item \textbf{Proof ($\impliedby$).}
Suppose $\M \cong B(\HH_R) \otimes \id_L$. Its center is $\mathbb C\id$ (as in the first double-commutant
example), so $\M$ is a factor. Take any one-dimensional projection $p = \ket{i}\bra{i}$ on $\HH_R$. Then
$P = p \otimes \id_L \in \M$. For any $A = a \otimes \id_L \in \M$:
\begin{align}
P A P
&\eqstep{1} (p a p) \otimes \id_L \notag\\
&\eqstep{2} \braket{i|a|i}\, (p \otimes \id_L)
\ \eqstep{3}\ \braket{i|a|i}\, P \ \in\ \mathbb{C} P . \notag
\end{align}
\stepwhy{1}{multiplying tensor products factor by factor, with $\id_L^3=\id_L$.}\quad
\stepwhy{2}{$pap=\ket i\braket{i|a|i}\bra i=\braket{i|a|i}\,p$.}\quad
\stepwhy{3}{the definition $P=p\otimes\id_L$.}

So every projection in $\M$ below $P$ is a multiple of $P$, hence $0$ or $P$. This means $P$ is minimal in
$\M$, and $\M$ is a type I factor. $\blacksquare$
\end{enumerate}
\end{keyresult}

So the statement ``the algebra describing this subsystem is type I'' is the precise, general form of the
statement ``the ordinary tensor-product picture of textbooks such as Griffiths and Sakurai applies here.''
Every worked example given so far in this chapter, including the two-qubit examples above, is type I by
construction. Types II and III, described next, are what becomes possible when this is no longer true.

\item \textbf{Type II: $\M$ has finite projections but no minimal ones.} This is a genuinely new possibility,
and it sounds strange at first. The finite/infinite distinction still works, and finite projections can be
compared in size using $d$. But there is no smallest nonzero size: below any nonzero projection there is
always a strictly smaller nonzero one. The values of $d$ are then forced to fill out a continuous range. Here
is the idea. In a type II factor, any finite projection $P$ can be split into two orthogonal pieces that are
equivalent to each other, $P=P_1+P_2$ with $P_1\sim P_2$ (stated without proof). Additivity then gives
$d(P_1)=d(P_2)=\tfrac12d(P)$. Repeating the split gives projections with $d=2^{-k}d(P)$ for every $k$, so $d$
takes arbitrarily small positive values. Adding such pieces gives every dyadic fraction of $d(P)$, and one can
show that every value in between is reached as well. So a von Neumann algebra can assign a \emph{real
number}, not only an integer, as the ``dimension'' of a subspace. Nothing like this happens in ordinary
finite-dimensional linear algebra, where dimension is a whole number obtained by counting basis vectors.
There are two subtypes, depending on whether $d(\id)$ is finite. In type $\mathrm{II}_1$ the identity has
finite dimension. We normalize $d(\id)=\tr(\id)=1$, so that $d(P)\in(0,1]$ for every nonzero projection: the
trace is normalized exactly like a probability. In type $\mathrm{II}_\infty$ we have $d(\id)=\infty$, so
$d(P)\in(0,\infty]$, and there is no natural normalization. Chapter~3 constructs a type $\mathrm{II}_1$ factor
by hand, from the $N\to\infty$ Bell-pair chain of Chapter~1. There, ``real-valued dimension'' stops being
abstract and can be watched in a specific family of projections.

\item \textbf{Type III: every nonzero projection is infinite.} This is the most extreme case. The
finite/infinite distinction has no content, because no projection except zero is finite. Formally,
$d(P)=\infty$ for every nonzero $P$. So there is no useful dimension function, and hence no (normal,
semifinite) trace on $\M$ at all. This looks like a mathematical dead end. How could one talk about
entanglement in an algebra with no density matrix, when density matrices are defined using a trace?
Chapter~4 meets exactly this obstacle and overcomes it with a different tool, Tomita--Takesaki modular
theory; that chapter is the technical heart of the notes. Type III is also the type that is physically
generic. The algebras of local regions in relativistic quantum field theory, and the boundary subalgebras of
holography in the strict large-$N$ limit, are both type III (in fact type $\mathrm{III}_1$, as later chapters
explain). This is not an exotic corner case. It is the everyday situation once you leave the world of
finitely many qubits.
\end{itemize}

Types II and III can sound, on first exposure, like unphysical mathematical curiosities. They are not. The
entangled spin chain of Chapter~1 already produces both, explicitly and by direct construction, and
Chapter~3 works through that construction completely.

\section{``Emergent'' Hilbert space and von Neumann algebra: the GNS construction}

\subsection*{Why this construction is needed}

Here is the gap again, now that every piece needed to close it is in hand. The completeness condition of a
von Neumann algebra (weak convergence) is stated in terms of matrix elements $\braket{\xi|A_n|\eta}$, so it
assumes that a Hilbert space already exists. The completeness condition of a $C^*$-algebra (norm
convergence) does not. It needs only an abstract norm with $\|A^\dagger A\|=\|A\|^2$. So you can specify a
$C^*$-algebra $\Alg$ abstractly, with no Hilbert space anywhere. Suppose you are also given a state $\omega$
on $\Alg$, in the sense defined earlier in this chapter: a positive, normalized linear functional, which needs
nothing beyond the algebra itself to define. Can a Hilbert space be \emph{built}, rather than assumed, in
which $\omega$ is represented by an ordinary vector? The Gelfand--Naimark--Segal (GNS) construction says yes,
always. It does so by an explicit, mechanical recipe.

\subsection*{The construction, built up in stages}

\textbf{Stage 1: treat algebra elements as candidate vectors.} This is the conceptual leap, so here it is in
words before the formulas. An element $C\in\Alg$ is reinterpreted as standing for ``the state you would get
by applying the operation $C$ to a fixed reference configuration.'' Two elements $C_1,C_2$ should count as
the \emph{same} state if no measurement, computed with $\omega$, could ever tell them apart.

\textbf{Stage 2: use $\omega$ to build an inner product on $\Alg$ itself.} For any $A,B\in\Alg$, define
\[
\braket{A|B} \equiv \omega(A^\dagger B) .
\]
Check that this deserves to be called an inner product. It is positive, $\braket{A|A}=\omega(A^\dagger A)\ge0$,
directly from the positivity of $\omega$. And it has the correct conjugate symmetry:
\begin{align}
\braket{B|A}=\omega(B^\dagger A)
&\eqstep{1} \omega\big((A^\dagger B)^\dagger\big) \notag\\
&\eqstep{2} \omega(A^\dagger B)^*
\ \eqstep{3}\ \braket{A|B}^* . \notag
\end{align}
\stepwhy{1}{$(A^\dagger B)^\dagger=B^\dagger A$.}\quad
\stepwhy{2}{$\omega(X^\dagger)=\omega(X)^*$, which holds for every positive functional (see the definitions
of weights and states).}\quad
\stepwhy{3}{the definition $\braket{A|B}\equiv\omega(A^\dagger B)$.}

Both required properties hold, using nothing but the definition of $\omega$.

\textbf{Stage 3: handle ``zero-length'' elements.} There may be nonzero $X\in\Alg$ with
$\omega(X^\dagger X)=0$. These are algebra elements that $\omega$ cannot see at all. Call the set of all such
$X$ the null space $\mathcal J$. The Cauchy--Schwarz inequality~\eqref{eq:CS-omega} shows that $\mathcal J$
is a linear subspace. A short computation shows more: if $X\in\mathcal J$ and $A\in\Alg$ is any element, then
$AX\in\mathcal J$ too. The proof is again an application of the Cauchy--Schwarz inequality:
\begin{align}
0 \ \le\ \omega\big((AX)^\dagger(AX)\big)
&\eqstep{1} \omega(X^\dagger A^\dagger AX) \notag \\
&\leqstep{2} \omega(X^\dagger X)^{1/2}\,\omega\big((A^\dagger AX)^\dagger A^\dagger AX\big)^{1/2}
\ \eqstep{3}\ 0 . \notag
\end{align}
\stepwhy{1}{regroup $(AX)^\dagger(AX)=X^\dagger A^\dagger AX$.}\quad
\stepwhy{2}{the Cauchy--Schwarz inequality for states,
$|\omega(A^\dagger B)|^2\le\omega(A^\dagger A)\,\omega(B^\dagger B)$, with $A$ replaced by $X$ and $B$
replaced by $A^\dagger AX$. This is the only inequality in the chain.}\quad
\stepwhy{3}{$\omega(X^\dagger X)=0$ because $X\in\mathcal J$, so the upper bound is exactly $0$ whatever the
second factor is.}

So $\omega\big((AX)^\dagger(AX)\big)=0$, which means $AX\in\mathcal J$. In the language of algebra,
$\mathcal J$ is a \emph{left ideal}. This computation does real work. It guarantees that quotienting out the
null elements in Stage~4 is consistent: multiplying by $A$ never turns a zero-length vector into a vector of
nonzero length.

\textbf{Stage 4: quotient and complete.} Declare $A$ and $A+X$ equivalent whenever $X\in\mathcal J$: algebra
elements that differ only by something $\omega$ cannot see count as the same vector. Let $[A]$ denote the
equivalence class of $A$. The inner product of Stage~2 does not depend on the choice of representative,
because $|\omega(X^\dagger B)|^2\le\omega(X^\dagger X)\,\omega(B^\dagger B)=0$ for $X\in\mathcal J$. On the
quotient space $\Alg/\mathcal J$ the inner product is positive definite: by construction, no nonzero vector
has zero length. Completing this space, by filling in the limits of Cauchy sequences just as the rational
numbers are completed to the real numbers, gives a Hilbert space, called $\HH_\omega$.

\textbf{Stage 5: define the representation.} The algebra acts on $\HH_\omega$ by left multiplication:
\[
\pi_\omega(A)\ket{[C]} \equiv \ket{[AC]} .
\]
This is well defined. If $C$ is replaced by $C+X$ with $X\in\mathcal J$, then $AC$ changes by $AX$, which is
also in $\mathcal J$ by Stage~3. It respects the multiplication of the algebra,
$\pi_\omega(A)\pi_\omega(B)=\pi_\omega(AB)$. Indeed, applying $\pi_\omega(B)$ and then $\pi_\omega(A)$ to $[C]$
gives $[A(BC)]=[(AB)C]$, which is $\pi_\omega(AB)$ applied to $[C]$. For a $C^*$-algebra, one also has
$\|\pi_\omega(A)[C]\|\le\|A\|\,\|[C]\|$, so each $\pi_\omega(A)$ is bounded and extends to the completion
$\HH_\omega$.

\textbf{Stage 6: identify the special vector.} The equivalence class of the identity element, $[\id]$, gets
its own name, $\ket\Omega\equiv\ket{[\id]}$. Two facts follow at once from the definitions above:
\[
\pi_\omega(A)\ket\Omega = \ket{[A\cdot\id]} = \ket{[A]}, \qquad\qquad
\omega(A) = \braket{\Omega|\pi_\omega(A)|\Omega} .
\]
The second follows from the first together with the definition of the inner product:
$\braket{\Omega|\pi_\omega(A)|\Omega} = \braket{[\id]|[A]} = \omega(\id^\dagger A)=\omega(A)$. So in the Hilbert
space just built, the vector $\ket\Omega$ reproduces the abstract state $\omega$ exactly, through the ordinary
formula $\braket{\Omega|A|\Omega}$ from undergraduate quantum mechanics. This is the purpose of the
construction. You gave it an abstract algebra and a number-valued rule $\omega$. It gave back a Hilbert space
and a vector inside it, and the usual formula for expectation values holds by construction, not by
assumption.

\begin{physconnect}[: Origin of the Vacuum $\ket0$]
$\ket\Omega$ is not found somewhere outside this construction. It \emph{is} the identity element of the
algebra, $\id\in\Alg$, renamed once the inner product has been put on $\Alg$ in Stage~2. A vector and the
Hilbert space it lives in are made in the same step, from the same two ingredients, $\Alg$ and $\omega$.
Neither exists before the other.

The construction also explains why $a$ annihilates $\ket\Omega$ for the oscillator ground state, without
using $\ket0$ as an input to define $\omega$. Start from one algebraic condition, $\omega(a^\dagger a)=0$,
together with positivity, $\omega(A^\dagger A)\ge0$ for every $A\in\Alg$, and normalization,
$\omega(\id)=1$. For any element $Y$, the Cauchy--Schwarz inequality~\eqref{eq:CS-omega} gives
$|\omega(Ya)|^2\le\omega(YY^\dagger)\,\omega(a^\dagger a)=0$. So $\omega(Ya)=0$ for every $Y$, and taking
adjoints, $\omega(a^\dagger Y)=\omega(Y^\dagger a)^*=0$ as well. In particular $\omega(a)=\omega(a^\dagger)=0$.
Using $[a,a^\dagger]=1$, every monomial in $a$ and $a^\dagger$ can be rewritten as a combination of
normal-ordered monomials $(a^\dagger)^ma^n$. Those with $m+n>0$ either end in $a$ or begin with $a^\dagger$,
so $\omega$ gives them zero. Only the constant term survives, and $\omega(\id)=1$ fixes its value. So the one
condition $\omega(a^\dagger a)=0$ determines $\omega$ completely. The oscillator box below uses the same
bookkeeping to show $\braket{0|a^n(a^\dagger)^n|0}=n!$. Here it is read as \emph{defining} $\omega$ from the
multiplication table of the algebra, not as evaluating a bra-ket that already existed.

Now apply Stage~3 to this $\omega$. The null ideal is $\mathcal J=\{X:\omega(X^\dagger X)=0\}$, and
$a\in\mathcal J$, because $\omega(a^\dagger a)=0$. Stage~4 quotients by exactly this $\mathcal J$, so every
element of $\mathcal J$ becomes the zero vector. So $[a]=0$, and $\pi_\omega(a)\ket\Omega=\ket{[a]}=0$. The
vacuum being annihilated by $a$ is not a separate fact about the world that $\omega$ happened to match. It is
what $[a]=0$ means, once $\mathcal J$ is fixed by the one number $\omega(a^\dagger a)$.

This is the general answer to where the vacuum and its Hilbert space come from. The vector space is not
searched for among candidates that already reproduce $\omega$. It \emph{is} $\Alg/\mathcal J$, completed. Its
points are exactly as fine as $\omega$ can distinguish, and no finer: two algebra elements become the same
vector exactly when $\omega$ assigns their difference zero length. Quotienting throws out the redundancy.
Completing, by filling in the limits of Cauchy sequences, turns this inner-product space into a Hilbert space
in the sense of Chapter~1.

In a field theory, $\omega$ is fixed in the same way, only more concretely, by a Euclidean path integral:
\[
\omega\big(O(x_1)\cdots O(x_n)\big) = \frac1Z\int\mathcal D\phi\;O(x_1)\cdots O(x_n)\,e^{-S[\phi]} .
\]
This formula computes every Euclidean correlator directly (Lorentzian correlators follow by analytic
continuation), and no Hilbert space, vector, or bra-ket appears in it. The CFT vacuum of Chapter~6 and the
JT-gravity states $\ket\beta$ of Chapter~9 are both defined this way. Canonical quantization, with its $\ket0$
and its Fock space, is the GNS representation built afterward from these correlators. Chapter~10 applies the
same idea to string theory. There, different asymptotic vacua are proposed to be different states $\omega$ on
one shared, background-independent algebra $\Alg_{\text{IIB}}$. Each has its own path-integral (or
correlator) definition that needs no Hilbert space. Each one's ``$\ket0$'' is simply $[\id]$ inside the GNS
space of \emph{that particular} $\omega$. It is never a single object shared across backgrounds. This is why
the resulting Hilbert spaces can be mutually inequivalent.
\end{physconnect}

The set $\{\pi_\omega(A)\ket\Omega : A\in\Alg\}$ is dense in $\HH_\omega$. This is automatic, because
$\HH_\omega$ was built as the completion of the set of classes $[A]=\pi_\omega(A)\ket\Omega$, so every vector
in $\HH_\omega$ can be approximated arbitrarily well by some $\pi_\omega(A)\ket\Omega$. A vector with this
property is called \textbf{cyclic}: every state can be reached from $\ket\Omega$ using only the algebra. This
is the same structure as the familiar construction of the harmonic-oscillator Hilbert space, by acting
repeatedly on the vacuum $\ket0$ with creation operators. The only difference is that here it is a theorem,
derived from the algebra and the state, rather than a separate postulate about how the Hilbert space is
built.

\begin{physconnect}[: GNS and Oscillator Fock Space]
The oscillator shows that GNS is the same calculation as the Fock-space construction you already know, run
in the opposite direction. We compute the same quantities in two ways and compare the answers.

\textbf{The ordinary way.} Start with a Hilbert space spanned by orthonormal vectors
$\ket0,\ket1,\ket2,\dots$, and define two operators by $a^\dagger\ket n=\sqrt{n+1}\,\ket{n+1}$ and
$a\ket n=\sqrt n\,\ket{n-1}$. Then $[a,a^\dagger]=1$, $a\ket0=0$, and $(a^\dagger)^m\ket0=\sqrt{m!}\,\ket m$.
Consequently
\[
\braket{0|a^n(a^\dagger)^m|0}=\sqrt{n!\,m!}\,\braket{n|m}=n!\,\delta_{nm} .
\]
Here the Hilbert space comes first, and the operators are defined on it.

\textbf{The GNS way.} Now forget the Hilbert space. Let $\Alg$ be the $*$-algebra generated by $1,a,a^\dagger$
with only the relation $[a,a^\dagger]=1$. Its elements are polynomials in $a$ and $a^\dagger$. Let $\omega$
be the state fixed by the single condition $\omega(a^\dagger a)=0$, as in the previous box. Since $a$ and
$a^\dagger$ are unbounded, $\Alg$ is a $*$-algebra but not a $C^*$-algebra. The GNS steps still go through,
with $\pi_\omega(a)$ and $\pi_\omega(a^\dagger)$ defined on the dense subspace of finite combinations of the
classes $[(a^\dagger)^n]$. Follow the recipe.

The Stage~2 inner product on $\Alg$ is $\braket{(a^\dagger)^n\,|\,(a^\dagger)^m}\equiv
\omega\big(a^n(a^\dagger)^m\big)$. First, it vanishes for $n\ne m$. Write $N=a^\dagger a$. The previous box
showed $\omega(a^\dagger Y)=0$ and $\omega(Ya)=0$ for every $Y$, so $\omega(NX)=\omega(XN)=0$ for every $X$.
Also $[N,a^n(a^\dagger)^m]=(m-n)\,a^n(a^\dagger)^m$, because $[N,a]=-a$ and $[N,a^\dagger]=a^\dagger$. Hence
\[
(m-n)\,\omega\big(a^n(a^\dagger)^m\big)=\omega(Na^n(a^\dagger)^m)-\omega(a^n(a^\dagger)^mN)=0 ,
\]
so $\omega\big(a^n(a^\dagger)^m\big)=0$ unless $n=m$. For $n=m$, write $\omega(X)$ as $\braket{0|X|0}$, as the
ordinary way would; the computation uses only the algebra and $\omega(Ya)=0$. It is one recursive step:
\begin{align}
\braket{0|a^n(a^\dagger)^n|0}
&\eqstep{1} \braket{0|a^{n-1}\big[a,(a^\dagger)^n\big]|0} + \braket{0|a^{n-1}(a^\dagger)^na|0} \notag\\
&\eqstep{2} n\,\braket{0|a^{n-1}(a^\dagger)^{n-1}|0} + 0 . \notag
\end{align}
\stepwhy{1}{split $a\,(a^\dagger)^n=[a,(a^\dagger)^n]+(a^\dagger)^n a$ and distribute.}\quad
\stepwhy{2}{the oscillator identity $[a,(a^\dagger)^n]=n(a^\dagger)^{n-1}$; the second term has the form
$\omega(Ya)$, which is zero.}

This is a recursion, $\braket{0|a^n(a^\dagger)^n|0}=n\,\braket{0|a^{n-1}(a^\dagger)^{n-1}|0}$, with base case
$\omega(\id)=1$ at $n=0$. Its solution is $n!$. (For example, $n=1$ gives $1\cdot1=1=1!$ and $n=2$ gives
$2\cdot1=2=2!$.) So the classes $[(a^\dagger)^n]$ are orthogonal, with squared norm $n!$. Normalizing gives
$\ket n\equiv[(a^\dagger)^n]/\sqrt{n!}$. The representation acts by $\pi_\omega(a^\dagger)[(a^\dagger)^n]=
[(a^\dagger)^{n+1}]$, which after normalization is $\pi_\omega(a^\dagger)\ket n=\sqrt{n+1}\,\ket{n+1}$. Also
$\pi_\omega(a)[(a^\dagger)^n]=[a(a^\dagger)^n]=n\,[(a^\dagger)^{n-1}]$ (System~3 of the next box shows this),
which after normalization is $\pi_\omega(a)\ket n=\sqrt n\,\ket{n-1}$. The cyclic vector $\ket\Omega=[\id]$ is
$\ket0$, and $\pi_\omega(a)\ket\Omega=[a]=0$.

\textbf{Comparison.} The two computations give the same answers:
\begin{center}
\small
\begin{tabular}{@{}lll@{}}
\toprule
& The ordinary way & The GNS way \\
\midrule
Inner products & $\braket{0|a^n(a^\dagger)^m|0}=n!\,\delta_{nm}$ & $\omega\big(a^n(a^\dagger)^m\big)=n!\,\delta_{nm}$ \\
Basis & $\ket n=(a^\dagger)^n\ket0/\sqrt{n!}$ & $\ket n=[(a^\dagger)^n]/\sqrt{n!}$ \\
Raising & $a^\dagger\ket n=\sqrt{n+1}\,\ket{n+1}$ & $\pi_\omega(a^\dagger)\ket n=\sqrt{n+1}\,\ket{n+1}$ \\
Lowering & $a\ket n=\sqrt n\,\ket{n-1}$ & $\pi_\omega(a)\ket n=\sqrt n\,\ket{n-1}$ \\
Ground state & $a\ket0=0$ & $\pi_\omega(a)\ket\Omega=[a]=0$ \\
\bottomrule
\end{tabular}
\end{center}
The map $\ket n\mapsto[(a^\dagger)^n]/\sqrt{n!}$ is a unitary between the two Hilbert spaces, and it carries
$a$ and $a^\dagger$ of the ordinary way to $\pi_\omega(a)$ and $\pi_\omega(a^\dagger)$. So the two routes give
the same Hilbert space, the same basis, the same matrix elements, and the same ground state. These are exact
equalities, not approximations.

Nothing here is new physics. It is the Fock-space construction from a first course on the harmonic
oscillator, run backward. What differs is only the direction of the logic. In the ordinary way you are handed
$\HH=L^2(\mathbb R)$ (or an abstract countable basis) first, and $a,a^\dagger$ are defined as operators on
it. In the GNS way the only inputs are the relation $[a,a^\dagger]=1$ and the state $\omega$ (equivalently,
``the vacuum is annihilated by $a$''). The whole Fock space, every $\ket n$ and every matrix element, is an
output. This is the sense in which the section on weights and states called the Hilbert space derived rather
than fundamental. For the oscillator nothing about the physics changes, because this $\omega$ produces the
usual Fock space. (For one degree of freedom this is no accident: by the Stone--von Neumann theorem, the
usual representation is, up to unitary equivalence, the only well-behaved irreducible one.) The cases that matter later
are those in which a different state produces an \emph{inequivalent} Hilbert space from the same
algebra. Examples are the $N$-Bell-pair chain, treated at the end of this chapter and in Chapters~3 and~4,
and the different asymptotic vacua of quantum gravity in Chapter~10.
\end{physconnect}

\begin{workedexamplebox}[: Concrete GNS Construction for Three Physical Systems]
To make the six-stage recipe fully mechanical, we now build the GNS Hilbert space, inner product, null ideal,
representation, and cyclic vector explicitly for three basic physical systems. In each case we also compare
the result with the answer ordinary quantum mechanics gives.

\begin{enumerate}[leftmargin=*]
\item \textbf{System 1: a pure state of a single qubit ($\Alg = M_2(\mathbb{C})$).}
\begin{itemize}[leftmargin=*]
\item \emph{Abstract algebra:} $\Alg = M_2(\mathbb{C})$, the 4-dimensional space of $2\times2$ complex
matrices, spanned by the matrix units $e_{ij} \equiv \ket{i}\bra{j}$ for $i,j \in \{0,1\}$.
\item \emph{State:} let $\omega_\uparrow$ be the pure spin-up state, $\omega_\uparrow(A) \equiv
\braket{0|A|0} = A_{00}$.
\item \emph{Pre-inner product on $\Alg$:} for any $A, B \in M_2(\mathbb{C})$,
\begin{align}
\braket{A|B}_\Alg = \omega_\uparrow(A^\dagger B)
&\eqstep{1} (A^\dagger B)_{00} \notag\\
&\eqstep{2} A_{00}^* B_{00} + A_{10}^* B_{10} . \notag
\end{align}
\stepwhy{1}{the definition $\omega_\uparrow(X)=X_{00}$.}\quad
\stepwhy{2}{matrix multiplication, $(A^\dagger B)_{00}=\sum_k (A^\dagger)_{0k}B_{k0}$, with $(A^\dagger)_{0k}=A_{k0}^*$.}

Only the first column of each matrix enters. The second-column entries $A_{01}, A_{11}, B_{01}, B_{11}$ do not
appear at all.
\item \emph{Null ideal $\mathcal{J}$:} an element $X \in \Alg$ is null exactly when
\[
\braket{X|X}_\Alg = |X_{00}|^2 + |X_{10}|^2 = 0 \iff X_{00} = 0 \quad \text{and} \quad X_{10} = 0 .
\]
So the null space is the 2-dimensional left ideal of matrices whose first column vanishes:
\[
\mathcal{J} = \left\{ \begin{pmatrix} 0 & a \\ 0 & b \end{pmatrix} : a,b \in \mathbb{C} \right\} = \mathrm{span}\{e_{01}, \, e_{11}\} .
\]
\item \emph{Quotient space $\HH_\omega = \Alg/\mathcal{J}$:} modulo $\mathcal{J}$, every matrix reduces to its
first column:
\begin{align}
[A] &\eqstep{1} A_{00} [e_{00}] + A_{10} [e_{10}] + A_{01}[e_{01}] + A_{11}[e_{11}] \notag\\
&\eqstep{2} A_{00}[e_{00}] + A_{10}[e_{10}] . \notag
\end{align}
\stepwhy{1}{write $A=\sum_{ij}A_{ij}e_{ij}$ and use linearity of $A\mapsto[A]$.}\quad
\stepwhy{2}{$e_{01},e_{11}\in\mathcal J$, so $[e_{01}]=[e_{11}]=0$.}

So the quotient space is 2-dimensional. Define the basis
\[
\ket{0} \equiv [e_{00}], \qquad \ket{1} \equiv [e_{10}] .
\]
Its inner products are $\braket{0|0} = 1$, $\braket{1|1} = 1$, $\braket{0|1} = 0$, so it is orthonormal, and
$\HH_\omega \cong \mathbb{C}^2$.
\item \emph{Representation $\pi_\omega$ and cyclic vector:} left multiplication by $A$ gives
\begin{align}
\pi_\omega(A)\ket{0} &= [A e_{00}] = A_{00}\ket{0} + A_{10}\ket{1}, \notag\\
\pi_\omega(A)\ket{1} &= [A e_{10}] = A_{01}\ket{0} + A_{11}\ket{1} . \notag
\end{align}
This is exactly ordinary matrix multiplication of $A$ on $\mathbb{C}^2$. The cyclic vector is the class of
the identity:
\begin{align}
\ket\Omega = [\id_2]
&\eqstep{1} [e_{00} + e_{11}] \notag\\
&\eqstep{2} [e_{00}]
\ \eqstep{3}\ \ket{0} . \notag
\end{align}
\stepwhy{1}{$\id_2=e_{00}+e_{11}$.}\quad
\stepwhy{2}{$e_{11}\in\mathcal J$, so it is zero in the quotient.}\quad
\stepwhy{3}{the definition $\ket0\equiv[e_{00}]$.}

\item \emph{Comparison with ordinary quantum mechanics:} a qubit in the state $\ket0$ is described by the
Hilbert space $\mathbb C^2$, the vector $\ket0$, and observables acting as $2\times2$ matrices. GNS gives
exactly this. It reduces the 4-dimensional matrix algebra to the familiar 2-dimensional qubit space
$\mathbb{C}^2$. Since $\omega_\uparrow$ is pure, the representation is irreducible: $\pi_\omega(\Alg)$ is all
of $B(\mathbb C^2)$, as the first structural fact after this box predicts.
\end{itemize}

\item \textbf{System 2: a mixed (thermal) state on $M_2(\mathbb{C})$, and the thermofield double.}
\begin{itemize}[leftmargin=*]
\item \emph{State:} let $\omega_\beta(A) \equiv \Tr(\rho_\beta A)$, with the full-rank density matrix
$\rho_\beta = \mathrm{diag}(p_0,p_1)$, where $p_0, p_1 > 0$ and $p_0 + p_1 = 1$.
\item \emph{Pre-inner product on $\Alg$:}
\begin{align}
\braket{A|B}_\Alg = \Tr(\rho_\beta A^\dagger B)
&= p_0 (A_{00}^* B_{00} + A_{10}^* B_{10}) \notag\\
&\quad + p_1 (A_{01}^* B_{01} + A_{11}^* B_{11}) . \notag
\end{align}
Each column of the matrix now enters, weighted by the corresponding eigenvalue of $\rho_\beta$.
\item \emph{Null ideal $\mathcal{J}$:} $\braket{X|X}_\Alg = 0$ means
\[
p_0 (|X_{00}|^2 + |X_{10}|^2) + p_1 (|X_{01}|^2 + |X_{11}|^2) = 0 .
\]
Because $p_0 > 0$ and $p_1 > 0$, every matrix entry must vanish, so $X = 0$. The null ideal is trivial,
$\mathcal{J} = \{0\}$.
\item \emph{The GNS Hilbert space:}
\[
\HH_\omega = \Alg / \{0\} \cong M_2(\mathbb{C}) \cong \mathbb{C}^2 \otimes \mathbb{C}^2 .
\]
The Hilbert space has dimension $4$, not $2$. Since $\|[e_{ij}]\|^2=\Tr(\rho_\beta e_{ji}e_{ij})=p_j$, the
normalized basis vectors are
\begin{align}
\ket{00} &\equiv \frac{1}{\sqrt{p_0}}[e_{00}], \qquad \ket{10} \equiv \frac{1}{\sqrt{p_0}}[e_{10}], \notag\\
\ket{01} &\equiv \frac{1}{\sqrt{p_1}}[e_{01}], \qquad \ket{11} \equiv \frac{1}{\sqrt{p_1}}[e_{11}] . \notag
\end{align}
The first label is the row index of the matrix and the second label is its column index.
\item \emph{The thermofield double vector:} the cyclic vector is
\[
\ket\Omega \equiv [\id_2] = [e_{00}] + [e_{11}] = \sqrt{p_0}\ket{00} + \sqrt{p_1}\ket{11} .
\]
This is an entangled pure state. When $p_j\propto e^{-\beta E_j}$ is a thermal distribution, it is exactly
the \textbf{thermofield double} (TFD) state. Left multiplication, $[Ae_{ij}]=\sum_kA_{ki}[e_{kj}]$, acts only
on the row label, so $\pi_\omega(\Alg) = M_2(\mathbb{C}) \otimes \id$. Its commutant is
$\pi_\omega(\Alg)' = \id \otimes M_2(\mathbb{C})$, acting on the column label.
\item \emph{Comparison with ordinary quantum mechanics:} the standard way to describe a mixed state $\rho$ by
a pure state is to purify it: add a second copy of the system and write
$\sum_j\sqrt{p_j}\,\ket j\ket j$. GNS produces exactly this purification. The second tensor factor, which
ordinary quantum mechanics introduces by hand, appears here automatically as the column index of the matrix.
\end{itemize}

\item \textbf{System 3: the bosonic oscillator (canonical commutation relation) and Fock space.}
\begin{itemize}[leftmargin=*]
\item \emph{Abstract algebra:} the $*$-algebra of polynomials in the ladder operators $a$ and $a^\dagger$,
with $[a, a^\dagger] = \id$.
\item \emph{State:} the vacuum state, defined algebraically by $\omega_0(\id) = 1$ and
$\omega_0(a^\dagger a) = 0$. Repeated use of $[a, a^\dagger] = \id$ fixes all the values
$\omega_0\big(a^m (a^\dagger)^n\big) = n!\,\delta_{mn}$, as shown in the box above.
\item \emph{Null ideal $\mathcal{J}$:} every element that ends with an annihilation operator is null, and the
null ideal is exactly $\mathcal{J} = \Alg a$. Modulo $\mathcal{J}$, every element reduces to a polynomial in
$a^\dagger$ alone: $[A] = \sum_n c_n [(a^\dagger)^n]$.
\item \emph{Inner product and orthonormal basis:}
\[
\Braket{[(a^\dagger)^m] \Big| [(a^\dagger)^n]} = \omega_0\big(a^m (a^\dagger)^n\big) = n!\,\delta_{mn} .
\]
Normalizing these vectors gives the standard Fock basis:
\[
\ket{n} \equiv \frac{1}{\sqrt{n!}}\big[(a^\dagger)^n\big], \qquad \braket{m|n} = \delta_{mn} .
\]
\item \emph{Action of the representation:}
\[
\pi_\omega(a^\dagger)\ket{n} = \sqrt{n+1}\,\ket{n+1}, \qquad \pi_\omega(a)\ket{n} = \sqrt{n}\,\ket{n-1} .
\]
The first follows from $\pi_\omega(a^\dagger)[(a^\dagger)^n]=[(a^\dagger)^{n+1}]$ and the normalization. For
the second,
\begin{align}
\big[a(a^\dagger)^n\big] = \big[[a,(a^\dagger)^n] + (a^\dagger)^n a\big]
&\eqstep{1} n\big[(a^\dagger)^{n-1}\big] + \big[(a^\dagger)^n a\big] \notag\\
&\eqstep{2} n\big[(a^\dagger)^{n-1}\big] . \notag
\end{align}
\stepwhy{1}{the commutator identity $[a,(a^\dagger)^n]=n(a^\dagger)^{n-1}$, the same one used in the box above.}\quad
\stepwhy{2}{$(a^\dagger)^na$ ends in an annihilation operator, so it lies in $\mathcal J=\Alg a$ and vanishes in
the quotient.}

Dividing by the normalizations gives $\pi_\omega(a)\ket n=n\sqrt{(n-1)!}/\sqrt{n!}\,\ket{n-1}=
\sqrt n\,\ket{n-1}$.
\item \emph{Comparison with ordinary quantum mechanics:} completing the space gives the infinite-dimensional
Fock space $\HH_{\text{Fock}} \cong \ell^2(\mathbb{N})$, with exactly the matrix elements of the ordinary
oscillator, as in the comparison table in the box above.
\end{itemize}
\end{enumerate}
\end{workedexamplebox}

Finally, the von Neumann algebra associated with all of this is the double commutant of the
representation,
\begin{equation}
\M\equiv\pi_\omega(\Alg)'' .
\label{eq:GNS-vN}
\end{equation}
This uses the double-commutant machinery from earlier in this chapter to promote the $C^*$-algebra
$\pi_\omega(\Alg)$, which is norm-closed but not necessarily weakly closed, to the von Neumann algebra it
generates.

\subsection*{Two structural facts, and what they mean}

The GNS triple $(\HH_\omega,\pi_\omega,\ket\Omega)$ is unique up to unitary equivalence. Any two
constructions that start from the same $\omega$ give the ``same'' Hilbert space, possibly described in a
different basis. We state this without proof. It is reassuring, because it means the construction does not
depend on any arbitrary choice made along the way.

Two further properties of $\omega$ translate directly into properties of the representation it produces:

\begin{itemize}
\item $\pi_\omega$ is \textbf{irreducible} if and only if $\omega$ is a \textbf{pure} state. Irreducible
means $\pi_\omega(\Alg)''=B(\HH_\omega)$: the representation generates \emph{every} bounded operator on
$\HH_\omega$, leaving nothing outside its reach.
\item If $\omega$ is \textbf{faithful}, then $\ket\Omega$ is \textbf{separating} for $\pi_\omega(\Alg)$.
Separating means that $\pi_\omega(A)\ket\Omega=0$ forces $\pi_\omega(A)=0$: no nonzero operator in the algebra
annihilates the reference vector. (Indeed, $\pi_\omega(A)\ket\Omega=[A]=0$ means $\omega(A^\dagger A)=0$, so
$A=0$.) Conversely, if $\ket\Omega$ is separating, then $\omega$ is faithful on the represented algebra
$\pi_\omega(\Alg)$.
\end{itemize}

The combination \emph{cyclic and separating} deserves attention now, although its full importance appears
only in Chapter~4. Cyclic says that the algebra reaches every vector, starting from $\ket\Omega$. Separating
says that different elements of the algebra act differently on $\ket\Omega$, so nothing nonzero is wasted.
The two notions are linked: a vector is cyclic for $\M$ if and only if it is separating for $\M'$ (a standard
fact). So a vector that is cyclic and separating for $\M$ is also cyclic and separating for $\M'$. In the
finite-dimensional examples of this chapter, this happens exactly when $\ket\Omega$ is entangled between the
two tensor factors with full Schmidt rank on both sides. This property is the foundation of Tomita--Takesaki
modular theory, the subject of Chapter~4.

\subsection*{The worked example, done completely with numbers}

We now work one standard example all the way through with numbers, rather than leaving it symbolic. It is
the most important computation in this chapter to see in full.

Let $\Alg=M_3(\mathbb C)$, the algebra of all $3\times3$ complex matrices, and take
\[
\rho = \begin{pmatrix}0.6&0&0\\0&0.4&0\\0&0&0\end{pmatrix} , \qquad \omega(A)\equiv\Tr(\rho A).
\]
First check that $\omega$ is a state in the sense defined earlier. It is linear, because the trace and matrix
multiplication are linear. It is positive: since $\rho$ is diagonal,
$\Tr(\rho A^\dagger A)=\sum_j\rho_{jj}(A^\dagger A)_{jj}=\sum_j\rho_{jj}\sum_k|A_{kj}|^2\ge0$. It is normalized:
$\Tr(\rho\,\id_3)=\Tr\rho=0.6+0.4+0=1$. The matrix $\rho$ has rank $2$, not $3$. It ignores the third basis
direction completely, so $\omega$ is \emph{not} faithful. For example, $e_{33}\ne0$ but
$\omega(e_{33}^\dagger e_{33})=\rho_{33}=0$. This property will show up directly in the construction.

Now follow the recipe, exactly as for System~2 above. The GNS inner product is
\begin{align}
\braket{A|B} = \Tr(\rho A^\dagger B)
&\eqstep{1} \sum_{j=1}^3 \rho_{jj}\,(A^\dagger B)_{jj} \notag\\
&\eqstep{2} 0.6\sum_{k=1}^3 A_{k1}^*B_{k1} + 0.4\sum_{k=1}^3 A_{k2}^*B_{k2} . \notag
\end{align}
\stepwhy{1}{$\rho$ is diagonal.}\quad
\stepwhy{2}{$(A^\dagger B)_{jj}=\sum_k A_{kj}^*B_{kj}$, and $\rho_{33}=0$ removes the $j=3$ term.}

So only the first two columns of a matrix matter. The null space $\mathcal J$ consists of the matrices whose
first two columns vanish, with the third column arbitrary. It is 3-dimensional, so the quotient
$\Alg/\mathcal J$ has dimension $9-3=6$. An orthonormal basis is
$\ket k_R\ket j_L\equiv[e_{kj}]/\sqrt{\rho_{jj}}$, for $k=1,2,3$ and $j=1,2$. Left multiplication,
$[Ae_{kj}]=\sum_lA_{lk}[e_{lj}]$, acts only on the row label $k$, so $\pi_\omega(A)=A\otimes\id_2$. The cyclic
vector is $[\id_3]=[e_{11}]+[e_{22}]+[e_{33}]=[e_{11}]+[e_{22}]$, since $e_{33}\in\mathcal J$. Altogether,
$\HH_\omega\cong\mathbb C^3\otimes\mathbb C^2$ is 6-dimensional: $3$ (the size of the matrices) times $2$ (the
rank of $\rho$). The cyclic vector and the representation are
\[
\ket\Omega=\sqrt{0.6}\,\ket1_R\ket1_L+\sqrt{0.4}\,\ket2_R\ket2_L ,
\qquad
\pi_\omega(A)=A\otimes\id_2 .
\]
We also checked this numerically. For several random complex $3\times3$ matrices $A,B$, we compared
$\omega(A)=\Tr(\rho A)$ with $\braket{\Omega|\pi_\omega(A)|\Omega}$, and the GNS inner product
$\braket{A|B}=\Tr(\rho A^\dagger B)$ with $\braket{\Omega|\pi_\omega(A)^\dagger\pi_\omega(B)|\Omega}$. In every
trial the two sides agreed to about one part in $10^{16}$, the rounding error of ordinary double-precision
arithmetic. The failure of faithfulness is also visible: $\pi_\omega(e_{13})\ket\Omega=0$ although
$e_{13}\ne0$. So $\ket\Omega$ is cyclic but not separating, as the second structural fact leads us to expect.

Recognize what this is. $\ket\Omega$ is the \textbf{purification} of $\rho$, familiar from quantum
information. Purification is the standard way to write a mixed state $\rho$ on a small Hilbert space as a
pure state on a larger one, by introducing an auxiliary system (here the $L$ factor) entangled with the
original. The GNS construction produces this purification automatically, with no separate step of ``now
introduce an auxiliary system.'' The auxiliary factor $\HH_L$ appears inside $\Alg/\mathcal J$ as the column
label of the matrix $[A]$. Its dimension is the rank of $\rho$, because $\mathcal J$ removes exactly the
columns on which $\rho$ vanishes.

Two variations of the same example show the two structural facts at work, on numbers you can check directly:
\begin{itemize}
\item Take $\rho=\mathrm{diag}(1,0,0)$ instead, which has rank $1$ and is a pure state. Then $\HH_\omega$
reduces to a single copy of $\mathbb C^3$: the $L$ factor becomes 1-dimensional and drops out. The
representation $\pi_\omega$ is irreducible, consistent with the first structural fact, since $\omega$ is now
pure.
\item Take $\rho=\tfrac13\id_3$ instead, which has full rank $3$ and is maximally mixed. Then
$\HH_\omega=\mathbb C^3\otimes\mathbb C^3$, and $\ket\Omega=\tfrac1{\sqrt3}\sum_{a=1}^3\ket a_R\ket a_L$ is both
cyclic \emph{and} separating. This is consistent with the second structural fact, since this $\omega$ is
faithful: $\omega(A^\dagger A)=\tfrac13\Tr(A^\dagger A)>0$ for every nonzero $A$. This maximally mixed state, with
its maximally entangled reference vector, is the $3$-dimensional analogue of the Bell pair at $\theta=\pi/4$
from Chapter~1. It is the same structure, one size larger.
\end{itemize}

\section{Emergent von Neumann algebras of the entangled spin example}

This closing section does one thing. It takes the GNS machinery just built and applies it explicitly to the
$N\to\infty$ Bell-pair chain of Chapter~1. This turns the informal description given there (``the
finite-energy excitations form some Hilbert space $\HH_{\Phi_\theta}$'') into an actual construction, with an
actual algebra and an actual state.

Define the algebra of \textbf{finite-energy operations}, $\Alg$. It is generated by operators of the form
\[
O = \alpha_1\otimes\alpha_2\otimes\cdots\otimes\alpha_n\otimes\cdots ,
\qquad \text{all but finitely many of the } \alpha_i \text{ equal to } \id_2\otimes\id_2 .
\]
So $O$ acts on only finitely many of the infinitely many spin pairs, and does nothing (the identity) to every
pair beyond some finite point. This restriction is not arbitrary. It is the algebraic form of the statement
that $O$ corresponds to a process reachable with finite energy. Recall from Chapter~1 that flipping
infinitely many spins costs infinite energy. An operator that acts nontrivially on infinitely many pairs at
once is exactly what a finite-energy process could never do. In the $N\to\infty$ limit, $O$ inherits a
well-defined norm from the finite-$N$ operators it is built from. Completing $\Alg$ in that norm makes it a
$C^*$-algebra. It is defined abstractly, with no particular Hilbert space singled out yet.

The reference state $\ket{\Phi_\theta}$, the infinite chain of Bell-like pairs from Chapter~1, defines a state
$\omega_\theta$ on this algebra in the obvious way:
\[
\omega_\theta(O) = \braket{\Phi_\theta|O|\Phi_\theta} .
\]
This makes sense because $O$ touches only finitely many pairs, so the expectation value uses only finitely
much of the state. No infinite sum or ill-defined limit is hidden in this definition.

Now apply the GNS construction of the previous section, exactly as given, to $(\Alg,\omega_\theta)$. It
produces a Hilbert space $\HH_{\Phi_\theta}$. This \emph{is} the space of finite-energy excitations described
informally in Chapter~1, now obtained as a mathematical construction rather than a descriptive placeholder.
Heuristically, and this is Stage~1 of the GNS recipe applied to this case, $\HH_{\Phi_\theta}$ is the completion
of the set of states reachable from $\ket{\Phi_\theta}$ by flipping a finite number of spins. That was
exactly the informal description in Chapter~1. In this representation, the von Neumann algebra you get
(the double commutant~\eqref{eq:GNS-vN}) is all of $B(\HH_{\Phi_\theta})$, every bounded operator on this
emergent Hilbert space. The reason is the first structural fact: $\ket{\Phi_\theta}$ is a pure state of the
whole chain, so the representation is irreducible.

Now suppose, as in Chapter~1, that you have access only to the $R$-half of the spin system. Let $\M_R$ be the
subalgebra of $B(\HH_{\Phi_\theta})$ consisting of operators that act only on the $R$-spins, completed under
weak convergence \emph{within this particular} $\HH_{\Phi_\theta}$. It is a von Neumann algebra. Its
\emph{type} (I, II, or III, depending on $\theta$) is the subject of Chapters~3 and~4. Define $\M_L$ in the
same way for the $L$-spins. Operations on the $L$-spins commute with operations on the $R$-spins, because they
act on physically distinct degrees of freedom. This shows directly that $\M_L\subseteq\M_R'$. In fact the two
are equal, $\M_L=\M_R'$. This equality is a theorem about infinite tensor products of this kind, which we
state without proof. So the $L$-algebra is exactly the commutant of the $R$-algebra. This is the
``complement of a subsystem'' picture from the start of this chapter, now realized concretely.

One further point removes an apparent asymmetry. This example is completely symmetric between $R$ and $L$:
nothing distinguishes them, and the construction could equally well have started from $\M_L$. It turns out
that the same Hilbert space $\HH_{\Phi_\theta}$ can be rebuilt using \emph{only} the $R$-algebra, with no
reference to $L$ at all. Concretely, redo the GNS construction using $\Alg_R$ (finite-energy operations on
the $R$-spins only) and the state $\omega_\theta(A)=\braket{\Phi_\theta|A|\Phi_\theta}$ restricted to
$A\in\Alg_R$. This restricted $\omega_\theta$ is a \emph{faithful} state of $\Alg_R$, because each pair's
reduced density matrix, $\mathrm{diag}(\cos^2\theta,\sin^2\theta)$, has full rank for $\theta\in(0,\pi/4]$. On
the full algebra $\Alg$, by contrast, $\omega_\theta$ is not faithful. For example, the projection
$\id-\ket{\phi_\theta}\bra{\phi_\theta}$ on a single pair is nonzero, but $\omega_\theta$ gives it zero. So
faithfulness is a property of a state \emph{relative to} a specific algebra, not a property of the state
alone. The GNS Hilbert space built from $(\Alg_R,\omega_\theta)$ is the same $\HH_{\Phi_\theta}$ as before,
because $\ket{\Phi_\theta}$ is cyclic for $\M_R$ (stated here without proof for the infinite chain).

The intuition is clearest for a single pair. Acting on $\ket{\phi_\theta}$ with operators on the right spin
alone already produces every vector in the pair's four-dimensional space $\mathbb C^2\otimes\mathbb C^2$.
Indeed,
\[
(A\otimes\id)\ket{\phi_\theta}=\cos\theta\,\big(A\ket0\big)\ket0+\sin\theta\,\big(A\ket1\big)\ket1 ,
\]
and as $A$ ranges over all $2\times2$ matrices, this covers all of $\mathbb C^2\otimes\mathbb C^2$ when
$\cos\theta$ and $\sin\theta$ are both nonzero. (We confirmed numerically that the four vectors
$(e_{ij}\otimes\id)\ket{\phi_\theta}$ are linearly independent.) So operators on the right spin reach
everything that operators on the left spin, or on both spins together, can reach. For the whole chain, the
same statement is that $\ket{\Phi_\theta}$ is cyclic for $\M_R$. This mirrors, on a much larger scale, the
point made earlier about minimal projections and irreducible representations. $\M_R$, as an abstract
algebra, already contains everything needed to rebuild the Hilbert space; $\HH_L$ only ever tracked
multiplicity. In this symmetric example, $R$ alone generates the same Hilbert space that $R$ and $L$ together
do.

With $\HH_{\Phi_\theta}$, $\M_R$, and $\M_L=\M_R'$ now built explicitly rather than only described, the stage is
set for the question that occupies the next two chapters. For different values of $\theta$, what
\emph{type}, in the precise sense of the classification above, is $\M_R$? Chapter~3 answers this by explicit
computation for $\theta=\pi/4$, where the answer is type $\mathrm{II}_1$. It also sets up the general case
$\theta\ne\pi/4$, which is type III and is completed in Chapter~4.
